% ============================================================
%  FLUX VIEWER ENGINE
%  Universal File Viewer & Editor — Complete System Design
%  Low-Level Android Native Architecture
%  Compiler: xelatex (run twice for TOC + refs)
% ============================================================
\documentclass[12pt,a4paper]{report}

% ── PACKAGES ────────────────────────────────────────────────
\usepackage{geometry}
\geometry{a4paper,top=24mm,bottom=24mm,left=20mm,right=20mm,headheight=15pt}
\usepackage{fontspec}
\usepackage{xcolor}
\usepackage{tikz}
\usetikzlibrary{shapes.geometric,arrows.meta,positioning,calc,fit,backgrounds}
\usepackage{tcolorbox}
\tcbuselibrary{skins,breakable,listings}
\usepackage{listings}
\usepackage{booktabs}
\usepackage{longtable}
\usepackage{array}
\usepackage{colortbl}
\usepackage{multirow}
\usepackage{tabularx}
\usepackage{fancyhdr}
\usepackage{titlesec}
\usepackage{enumitem}
\usepackage{hyperref}
\usepackage{mdframed}
\usepackage{amsmath}
\usepackage{amssymb}
\usepackage{microtype}
\usepackage{parskip}
\usepackage{setspace}
\usepackage{caption}
\usepackage{float}
\usepackage{makecell}
\usepackage{ragged2e}
\usepackage{adjustbox}

% ── FONTS ───────────────────────────────────────────────────
\setmainfont{DejaVu Serif}
\setsansfont{DejaVu Sans}
\setmonofont{DejaVu Sans Mono}[Scale=0.86]

% ── COLOURS ─────────────────────────────────────────────────
\definecolor{P1}  {HTML}{0D47A1}  % deep blue    — primary
\definecolor{P2}  {HTML}{1565C0}  % indigo
\definecolor{P3}  {HTML}{00838F}  % cyan
\definecolor{P4}  {HTML}{00695C}  % teal
\definecolor{G1}  {HTML}{1B5E20}  % deep green
\definecolor{OR}  {HTML}{E65100}  % orange
\definecolor{RD}  {HTML}{B71C1C}  % red
\definecolor{PU}  {HTML}{4A148C}  % purple
\definecolor{GY}  {HTML}{37474F}  % gray
\definecolor{LGY} {HTML}{607D8B}  % light gray
\definecolor{BRN} {HTML}{4E342E}  % brown
\definecolor{TEAL}{HTML}{006064}  % dark teal
% fills
\definecolor{fB}  {HTML}{E3F2FD}
\definecolor{fC}  {HTML}{E0F7FA}
\definecolor{fG}  {HTML}{E8F5E9}
\definecolor{fO}  {HTML}{FFF3E0}
\definecolor{fR}  {HTML}{FFEBEE}
\definecolor{fP}  {HTML}{F3E5F5}
\definecolor{fY}  {HTML}{FFFDE7}
\definecolor{fN}  {HTML}{E8EAF6}
\definecolor{fT}  {HTML}{E0F2F1}
\definecolor{cBG} {HTML}{1A1A2E}
\definecolor{cFG} {HTML}{A8D8A8}
\definecolor{cKW} {HTML}{80CBC4}
\definecolor{cST} {HTML}{FFD180}
\definecolor{cCM} {HTML}{78909C}
\definecolor{rowA}{HTML}{E3F2FD}
\definecolor{rowB}{HTML}{FFFFFF}

% ── HYPERREF ────────────────────────────────────────────────
\hypersetup{colorlinks=true,linkcolor=P1,urlcolor=P3,citecolor=P4,
  pdftitle={FLUX Viewer Engine — Universal File System Design},
  bookmarksnumbered=true,bookmarksopen=true}

% ── SECTION STYLES ──────────────────────────────────────────
\titleformat{\chapter}[block]
  {\color{P1}\sffamily\bfseries\LARGE}
  {\colorbox{P1}{\color{white}\enspace CHAPTER \thechapter\enspace}}{1em}{}
  [\vspace{2pt}\color{P3}\rule{\linewidth}{2pt}]
\titleformat{\section}
  {\color{P1}\sffamily\bfseries\Large}{\thesection}{0.8em}{}
  [\color{P3}\rule{\linewidth}{0.8pt}]
\titleformat{\subsection}
  {\color{P3}\sffamily\bfseries\large}{\thesubsection}{0.8em}{}
\titleformat{\subsubsection}
  {\color{GY}\sffamily\bfseries\normalsize}{\thesubsubsection}{0.6em}{}
\titlespacing{\chapter}   {0pt}{0pt}{14pt}
\titlespacing{\section}   {0pt}{16pt}{7pt}
\titlespacing{\subsection}{0pt}{12pt}{5pt}

% ── PAGE STYLE ──────────────────────────────────────────────
\pagestyle{fancy}\fancyhf{}
\fancyhead[L]{\sffamily\small\color{LGY}FLUX Viewer Engine}
\fancyhead[R]{\sffamily\small\color{LGY}\leftmark}
\fancyfoot[C]{\sffamily\small\color{LGY}\thepage}
\renewcommand{\headrulewidth}{0.4pt}
\renewcommand{\headrule}{\color{P1}\hrule width\headwidth height\headrulewidth}

% ── LISTINGS ────────────────────────────────────────────────
\lstdefinestyle{flux}{
  backgroundcolor=\color{cBG},basicstyle=\ttfamily\small\color{cFG},
  keywordstyle=\color{cKW}\bfseries,stringstyle=\color{cST},
  commentstyle=\color{cCM}\itshape,numberstyle=\tiny\color{LGY},
  numbers=left,numbersep=8pt,frame=none,breaklines=true,
  keepspaces=true,showstringspaces=false,tabsize=2,
  xleftmargin=10pt,xrightmargin=4pt,aboveskip=6pt,belowskip=6pt,
  columns=flexible,
  morekeywords={fun,val,var,class,data,object,override,companion,suspend,
    when,is,in,return,if,else,for,while,true,false,null,private,public,
    internal,sealed,enum,interface,abstract,import,package,by,lazy,init,
    constructor,it,this,super,async,await,const,final,static,void,String,
    Int,Long,Boolean,Double,Float,Short,Byte,Unit,Any,List,Map,Set,
    IntArray,ByteArray,FloatArray,MappedByteBuffer,FileChannel,Canvas,
    Paint,Rect,RectF,Bitmap,Path,Matrix,PorterDuffXfermode},
}
\lstset{style=flux}

% ── TCOLORBOXES ─────────────────────────────────────────────
\newtcolorbox{ibox}[1][]{enhanced,breakable,colback=fB,colframe=P1,
  leftrule=4pt,rightrule=0.4pt,toprule=0.4pt,bottomrule=0.4pt,arc=2pt,
  boxsep=4pt,fonttitle=\sffamily\bfseries\color{P1},title={#1}}
\newtcolorbox{wbox}[1][]{enhanced,breakable,colback=fO,colframe=OR,
  leftrule=4pt,rightrule=0.4pt,toprule=0.4pt,bottomrule=0.4pt,arc=2pt,
  boxsep=4pt,fonttitle=\sffamily\bfseries\color{OR},title={#1}}
\newtcolorbox{cbox}[1][]{enhanced,breakable,colback=fR,colframe=RD,
  leftrule=4pt,rightrule=0.4pt,toprule=0.4pt,bottomrule=0.4pt,arc=2pt,
  boxsep=4pt,fonttitle=\sffamily\bfseries\color{RD},title={#1}}
\newtcolorbox{gbox}[1][]{enhanced,breakable,colback=fG,colframe=P4,
  leftrule=4pt,rightrule=0.4pt,toprule=0.4pt,bottomrule=0.4pt,arc=2pt,
  boxsep=4pt,fonttitle=\sffamily\bfseries\color{P4},title={#1}}
\newtcolorbox{pbox}[1][]{enhanced,breakable,colback=fP,colframe=PU,
  leftrule=4pt,rightrule=0.4pt,toprule=0.4pt,bottomrule=0.4pt,arc=2pt,
  boxsep=4pt,fonttitle=\sffamily\bfseries\color{PU},title={#1}}

% ── TABLE HELPERS ───────────────────────────────────────────
\newcommand{\ra}[1]{\renewcommand{\arraystretch}{#1}}
\newcolumntype{L}[1]{>{\RaggedRight\arraybackslash}p{#1}}
\newcolumntype{C}[1]{>{\centering\arraybackslash}p{#1}}

\newcommand{\TH}[1]{\rowcolor{P1}\color{white}\bfseries #1}
\newcommand{\THC}[1]{\rowcolor{P3}\color{white}\bfseries #1}
\newcommand{\THG}[1]{\rowcolor{GY}\color{white}\bfseries #1}
\newcommand{\THO}[1]{\rowcolor{OR}\color{white}\bfseries #1}
\newcommand{\THRD}[1]{\rowcolor{RD}\color{white}\bfseries #1}
\newcommand{\THPU}[1]{\rowcolor{PU}\color{white}\bfseries #1}
\newcommand{\THGN}[1]{\rowcolor{G1}\color{white}\bfseries #1}

\newcommand{\PB}{\newpage}

% ── SECTION PART BANNER ─────────────────────────────────────
\newcommand{\partbanner}[2]{%
\begin{tcolorbox}[enhanced,colback=#2,colframe=#2,arc=3pt,
  boxsep=8pt,left=10pt,right=10pt]
  \sffamily\bfseries\LARGE\color{white}#1
\end{tcolorbox}\vspace{4pt}}

% ============================================================
\begin{document}
% ============================================================

% ── TITLE PAGE ──────────────────────────────────────────────
\begin{titlepage}
\pagecolor{P1}\color{white}\centering
\vspace*{20mm}
{\sffamily\fontsize{56}{62}\selectfont\bfseries FLUX VIEWER ENGINE}
\vspace{6mm}

\color{white!70!P1}\rule{0.65\linewidth}{1.5pt}\vspace{5mm}

\color{white}{\sffamily\Large\bfseries Universal File Viewer \& Editor}\\[3mm]
{\sffamily\large Complete Low-Level System Design for Android}

\vspace{10mm}
{\sffamily\normalsize
  PDF · DOCX · XLSX · PPTX · Images · Video · Audio\\[3pt]
  TXT · HTML · Code · SQL · CSV · JSON · XML · YAML\\[3pt]
  ZIP · APK · SVG · EPUB · Markdown · Font Files
}

\vspace{12mm}\color{white!80!P1}\rule{0.45\linewidth}{0.6pt}\vspace{6mm}

{\sffamily\small
  Zero Heavy Libraries \enspace·\enspace Native Kotlin Only\\[3pt]
  Memory-Mapped I/O \enspace·\enspace Tile-Based Rendering\\[3pt]
  Progressive Loading \enspace·\enspace Sub-100ms First Paint
}

\vfill
{\sffamily\footnotesize\color{white!55!P1}
  FLUX Engineering \enspace·\enspace Version 1.0 \enspace·\enspace 2025 \enspace·\enspace Confidential
}
\end{titlepage}
\pagecolor{white}
\tableofcontents
\newpage

% ============================================================
% CHAPTER 1 — MARKET ANALYSIS
% ============================================================
\chapter{Market Analysis — Why Every Current File Viewer Is Broken}

\section{The Fundamental Problem}

Every existing file viewer on Android is slow for the same root reasons. Before
designing a solution, we must diagnose the disease precisely — not just the
symptom.

\begin{cbox}[The Three Root Diseases of File Viewing on Android]
\begin{enumerate}[leftmargin=1.4em,topsep=2pt,itemsep=5pt]
\item \textbf{Separate-app architecture:} Every file type requires launching a
  different app. Cold start of WPS Office: 3--5\,s. Adobe Acrobat: 2--4\,s.
  Google Drive PDF viewer: 2--3\,s. The user waits before seeing a single byte.
\item \textbf{Monolithic loading:} Current apps load the \emph{entire} file into
  memory before showing anything. A 50-page PDF is fully parsed + rasterised
  before page 1 appears.
\item \textbf{Heavyweight library dependency:} PDFium (Chromium's PDF engine) is
  38\,MB. LibreOffice Mobile is 50\,MB+. These libraries are designed for
  desktop-class hardware and carry massive initialisation overhead.
\end{enumerate}
\end{cbox}

\section{Detailed Analysis — App by App}

\subsection{PDF Viewers}

\begin{center}\ra{1.35}
\begin{tabular}{L{3.0cm} L{2.8cm} L{2.8cm} L{5.2cm}}
\TH{App} & \TH{Library Used} & \TH{Cold Start} & \TH{Root Cause of Slowness} \\
\rowcolor{rowA} Adobe Acrobat & Proprietary C++ engine (Acrobat SDK) & 3--5\,s &
  Loads entire Acrobat font renderer, 80MB+ RAM on launch \\
\rowcolor{rowB} Google Drive PDF & PDFium (Chromium fork, 38MB) & 2--3\,s &
  Full PDFium init even for 1-page PDF; renders all pages to off-screen bitmap \\
\rowcolor{rowA} MuPDF & MuPDF C library (via JNI) & 1.5--2.5\,s &
  JNI bridge overhead; full document object model built before display \\
\rowcolor{rowB} Xodo & PDFium + custom JS engine & 4--6\,s &
  JS engine loads for annotation/form support even on read-only PDFs \\
\rowcolor{rowA} Built-in (Samsung/Xiaomi) & PDFium variant & 2--4\,s &
  Same PDFium root cause; no tile rendering \\
\end{tabular}
\end{center}

\textbf{The Real Problem With PDF:} Android has had
\texttt{android.graphics.pdf.PdfRenderer} built into the OS since API 21
(Android 5.0). It is hardware-accelerated, runs at native speed, requires zero
external library, and is available on 99.9\% of Android devices. \emph{Nobody
uses it for a proper file viewer because they wrap it naively} — rendering the
whole page at full resolution before showing anything.

\subsection{Office Document Viewers (DOCX, XLSX, PPTX)}

\begin{center}\ra{1.35}
\begin{tabular}{L{3.0cm} L{3.0cm} L{2.4cm} L{5.4cm}}
\TH{App} & \TH{Approach} & \TH{Open Time} & \TH{Root Cause} \\
\rowcolor{rowA} WPS Office & Full LibreOffice-style layout engine & 3--6\,s &
  Full font substitution + style computation before page 1 \\
\rowcolor{rowB} Microsoft Word (Android) & Cloud-assisted render + local OOXML parser & 2--5\,s &
  Uploads file to Microsoft servers for rendering if complex; full XML parse locally \\
\rowcolor{rowA} Google Docs (offline) & Converts to internal format on open & 4--8\,s &
  Full format conversion: OOXML → Google's internal binary format on first open \\
\rowcolor{rowB} Polaris Office & Proprietary OOXML engine (C++ via JNI) & 2--4\,s &
  JNI init overhead; builds complete page layout model before any rendering \\
\rowcolor{rowA} QuickEdit & Text-only, no rich formatting & 0.3--0.8\,s &
  Reads file as plain text — works but loses all formatting \\
\end{tabular}
\end{center}

\begin{ibox}[The OOXML Secret]
DOCX, XLSX, and PPTX are ZIP archives containing XML files. There is
\emph{nothing proprietary} about them — they are an ISO standard (ISO/IEC 29500).
A DOCX file with 20 pages of basic text and tables uses fewer than 15 XML
elements from the OOXML specification. Full office suites parse 500+
element types because they support 100\% of the spec. We only need to
parse what 95\% of documents actually use.
\end{ibox}

\subsection{Image Viewers}

\begin{center}\ra{1.35}
\begin{tabular}{L{3.2cm} L{3.4cm} L{2.0cm} L{5.2cm}}
\TH{App} & \TH{Approach} & \TH{Problem} & \TH{Why It's Slow} \\
\rowcolor{rowA} Google Photos & Full bitmap decode to RAM & OOM on large files &
  \texttt{BitmapFactory.decodeFile()} decodes entire 36MB RAW into 36MB+ RAM \\
\rowcolor{rowB} Samsung Gallery & Same: full decode & OOM possible &
  No \texttt{BitmapRegionDecoder} for progressive tile loading \\
\rowcolor{rowA} Quickpic / Piktures & Full decode + Glide library & 200--500ms per image &
  Glide adds disk cache + memory cache overhead even for local files \\
\rowcolor{rowB} Files by Google & Thumbnail only, no full viewer & Cannot zoom & Opens external app for full view \\
\end{tabular}
\end{center}

\textbf{The Fix:} Android has had \texttt{BitmapRegionDecoder} since API 10
(Android 2.3). It decodes \emph{only the visible rectangle} of an image. A
100\,MP photo displayed at 360dp width needs only $360 \times 360 \times 4$\,bytes
= 500\,KB decoded, not 400\,MB.

\subsection{Code / Text / SQL / HTML Viewers}

\begin{center}\ra{1.35}
\begin{tabular}{L{3.0cm} L{3.2cm} L{5.6cm}}
\TH{File Type} & \TH{Current Solution} & \TH{Problem} \\
\rowcolor{rowA} SQL files & No dedicated viewer & Must install DBeaver Mobile or open as text \\
\rowcolor{rowB} HTML files & WebView or browser & WebView cold start 500ms+; browser requires internet context \\
\rowcolor{rowA} Source code (.kt/.py/.js) & Text editor apps (QuickEdit, Acode) &
  Load entire file into \texttt{TextView}; 10MB source file → 10MB heap allocation, ANR \\
\rowcolor{rowB} JSON/XML & Online tools / text editor & No tree view; massive files cause OOM in TextView \\
\rowcolor{rowA} CSV & Excel / WPS & Full spreadsheet engine for what is basically delimited text \\
\rowcolor{rowB} Markdown & Typora / Obsidian (heavy) & Full Electron/WebView-based apps; 200MB+ RAM \\
\end{tabular}
\end{center}

\section{The Performance Gap Summary}

\begin{center}\ra{1.4}
\begin{tabular}{L{3.0cm} C{2.4cm} C{2.4cm} C{2.4cm} C{3.6cm}}
\THRD{File Type} & \THRD{Best Current App} &
\THRD{Current Speed} & \THRD{FLUX Target} & \THRD{Our Advantage} \\
\rowcolor{rowA} PDF (10 pages) & MuPDF & 1.5\,s cold start & \textbf{$<$80ms} & PdfRenderer + tile cache \\
\rowcolor{rowB} DOCX (20 pages) & WPS Office & 3\,s cold start & \textbf{$<$150ms} & Minimal OOXML parser \\
\rowcolor{rowA} XLSX (50 rows) & WPS Office & 2\,s & \textbf{$<$100ms} & Virtual grid renderer \\
\rowcolor{rowB} Image (12MP) & Google Photos & 400\,ms full decode & \textbf{$<$30ms} & BitmapRegionDecoder tile \\
\rowcolor{rowA} 10MB text file & QuickEdit & ANR / slow scroll & \textbf{$<$50ms} & Line-virtual Canvas renderer \\
\rowcolor{rowB} HTML file & WebView & 500ms+ & \textbf{$<$200ms} & System WebView + no cold start \\
\rowcolor{rowA} JSON (5MB) & Text editor & OOM / slow & \textbf{$<$50ms} & Streaming tree parser \\
\rowcolor{rowB} ZIP file & ES File Explorer & 1\,s & \textbf{$<$30ms} & \texttt{ZipInputStream} native \\
\rowcolor{rowA} SQL file & None & No viewer & \textbf{$<$30ms} & Syntax-highlighted text renderer \\
\rowcolor{rowB} SVG & Browser & 800ms & \textbf{$<$100ms} & Native Canvas SVG renderer \\
\end{tabular}
\end{center}

\PB

% ============================================================
% CHAPTER 2 — FLUX VIEWER ENGINE ARCHITECTURE
% ============================================================
\chapter{FLUX Viewer Engine — Core Architecture}

\section{Design Philosophy}

\begin{pbox}[The Five Laws of FLUX Viewer Engine]
\begin{enumerate}[leftmargin=1.4em,topsep=2pt,itemsep=5pt]
\item \textbf{First paint in $<$100ms.} Show something — anything — before
  100ms. Even a placeholder tile is better than a blank screen.
\item \textbf{Never load the full file.} Always use streaming, tiling, or
  virtual rendering. A 500MB video, a 100-page PDF, a 10MB JSON file — none
  is ever fully in RAM at once.
\item \textbf{Zero external libraries for core rendering.} Every rendering
  primitive uses Android's built-in APIs: \texttt{PdfRenderer}, \texttt{Canvas},
  \texttt{BitmapRegionDecoder}, \texttt{MediaPlayer}, \texttt{WebView}.
\item \textbf{Hardware-accelerated always.} All drawing via \texttt{Canvas} on a
  hardware-accelerated \texttt{View}. Never \texttt{setLayerType(SOFTWARE)}.
\item \textbf{Edit layer is separate from the render layer.} The viewer renders
  the original file bytes. Edits are an in-memory overlay written back only on
  explicit save. This keeps the viewer always fast and edits always safe.
\end{enumerate}
\end{pbox}

\section{System Architecture — Three Layers}

\begin{center}\ra{1.4}
\begin{tabular}{C{0.8cm} L{2.6cm} L{3.2cm} L{4.0cm} L{4.2cm}}
\TH{Layer} & \TH{Name} & \TH{Language} & \TH{Responsibility} & \TH{Key Classes} \\
\rowcolor{rowA} L1 & Source Layer & Kotlin & Memory-mapped file reading, format detection, streaming parsers &
  \texttt{MmapSource}, \texttt{FormatDetector}, \texttt{StreamParser} \\
\rowcolor{rowB} L2 & Render Layer & Kotlin + Canvas & Tile generation, viewport management, hardware-accelerated drawing &
  \texttt{TileEngine}, \texttt{ViewportManager}, \texttt{FluxCanvasView} \\
\rowcolor{rowA} L3 & Edit Layer & Kotlin & In-memory edit overlay, diff tracking, write-back on save &
  \texttt{EditOverlay}, \texttt{DiffTracker}, \texttt{FileSaveEngine} \\
\end{tabular}
\end{center}

\section{The Core: MmapSource — Memory-Mapped File Reading}

Every format parser in FLUX reads through a single abstraction:
\texttt{MmapSource}. It memory-maps the file into the process's virtual address
space. The OS loads pages on demand — we never read the whole file unless the
user scrolls through the whole file.

\begin{lstlisting}[language=Kotlin,caption={MmapSource — zero-copy file access}]
class MmapSource(private val file: File) : AutoCloseable {
  private val channel: FileChannel =
    FileInputStream(file).channel

  // Memory-map entire file — OS loads pages on demand
  private val buffer: MappedByteBuffer =
    channel.map(FileChannel.MapMode.READ_ONLY, 0, file.length())

  val size: Long = file.length()

  // Read bytes at absolute offset — O(1), zero-copy
  fun readBytes(offset: Long, length: Int): ByteArray {
    val buf = ByteArray(length)
    buffer.position(offset.toInt())
    buffer.get(buf)
    return buf
  }

  // Read a single byte — O(1)
  fun readByte(offset: Long): Byte = buffer.get(offset.toInt())

  // Read little-endian Int32 at offset — O(1)
  fun readInt32LE(offset: Long): Int {
    val b = buffer
    return (b.get(offset.toInt()).toInt() and 0xFF) or
           ((b.get((offset+1).toInt()).toInt() and 0xFF) shl 8) or
           ((b.get((offset+2).toInt()).toInt() and 0xFF) shl 16) or
           ((b.get((offset+3).toInt()).toInt() and 0xFF) shl 24)
  }

  // Slice a region as a new view — O(1), no copy
  fun slice(offset: Long, length: Int): MmapSlice {
    buffer.position(offset.toInt())
    val sliced = buffer.slice()
    sliced.limit(length)
    return MmapSlice(sliced)
  }

  override fun close() { channel.close() }
}
\end{lstlisting}

\section{The Tile Engine — Universal Rendering Primitive}

Every format — PDF pages, spreadsheet grids, long text files, large images —
is decomposed into \emph{tiles}. A tile is a fixed-size rectangle of rendered
content. Only tiles intersecting the current viewport are rendered.

\begin{lstlisting}[language=Kotlin,caption={TileEngine — format-agnostic tile management}]
data class TileKey(val row: Int, val col: Int, val zoomLevel: Int)
data class Tile(val key: TileKey, val bitmap: Bitmap, val renderedAt: Long)

class TileEngine(
  private val tileW: Int = 512,   // tile width in px
  private val tileH: Int = 512,   // tile height in px
  private val maxCacheMB: Int = 40 // max tile cache RAM
) {
  // LRU cache: evict oldest tile when RAM limit hit
  private val cache = object : LruCache<TileKey, Tile>(maxCacheMB * 1024 * 1024) {
    override fun sizeOf(key: TileKey, value: Tile): Int =
      value.bitmap.byteCount
  }

  // Request a tile — returns cached or triggers async render
  fun getTile(key: TileKey, renderer: TileRenderer): Tile? {
    cache.get(key)?.let { return it }  // cache hit O(1)

    // Cache miss: submit to render queue (background thread)
    renderQueue.submit {
      val bitmap = renderer.renderTile(key, tileW, tileH)
      val tile   = Tile(key, bitmap, System.currentTimeMillis())
      cache.put(key, tile)
      onTileReady(tile)  // callback to invalidate view region
    }
    return null  // caller shows placeholder
  }

  // Which tiles intersect the viewport?
  fun tilesForViewport(viewport: RectF, contentSize: SizeF): List<TileKey> {
    val startCol = (viewport.left / tileW).toInt().coerceAtLeast(0)
    val endCol   = (viewport.right / tileW).toInt()
    val startRow = (viewport.top / tileH).toInt().coerceAtLeast(0)
    val endRow   = (viewport.bottom / tileH).toInt()
    return (startRow..endRow).flatMap { row ->
      (startCol..endCol).map { col -> TileKey(row, col, currentZoom) }
    }
  }
}

// Each format implements this interface
interface TileRenderer {
  fun renderTile(key: TileKey, tileW: Int, tileH: Int): Bitmap
}
\end{lstlisting}

\section{Format Detection — Before Parsing}

FLUX detects file type from the first 16 bytes (magic bytes), not from
the extension. Extensions lie; magic bytes don't.

\begin{lstlisting}[language=Kotlin,caption={FormatDetector — magic byte identification}]
object FormatDetector {
  fun detect(source: MmapSource): FileFormat {
    if (source.size < 4) return FileFormat.BINARY

    val magic = source.readBytes(0, 16)
    return when {
      // PDF: starts with '%PDF-'
      magic.startsWith("%PDF-")     -> FileFormat.PDF
      // ZIP-based: DOCX, XLSX, PPTX, EPUB, APK all start with PK (0x50 0x4B)
      magic[0] == 0x50.b && magic[1] == 0x4B.b -> detectZipBased(source)
      // JPEG: FF D8 FF
      magic[0] == 0xFF.b && magic[1] == 0xD8.b -> FileFormat.JPEG
      // PNG: 89 50 4E 47
      magic.startsWith("\u0089PNG")  -> FileFormat.PNG
      // WebP: RIFF....WEBP
      magic.startsWith("RIFF") && magic.sliceArray(8..11).toString(Charsets.US_ASCII) == "WEBP"
                                    -> FileFormat.WEBP
      // GIF: GIF87a or GIF89a
      magic.startsWith("GIF8")      -> FileFormat.GIF
      // HEIC/HEIF: ftyp box at offset 4
      source.readInt32LE(4) == 0x66747970  -> FileFormat.HEIC
      // MP4/MOV: ftyp box
      magic.sliceArray(4..7).toString(Charsets.US_ASCII) in listOf("ftyp","moov")
                                    -> FileFormat.VIDEO_MP4
      // MP3: ID3 header or 0xFF 0xFB
      magic.startsWith("ID3") || (magic[0] == 0xFF.b && (magic[1].toInt() and 0xE0) == 0xE0)
                                    -> FileFormat.AUDIO_MP3
      // FLAC: fLaC
      magic.startsWith("fLaC")      -> FileFormat.AUDIO_FLAC
      // SQLite: starts with 'SQLite format 3'
      magic.startsWith("SQLite")    -> FileFormat.SQLITE
      // Text heuristic: all printable ASCII in first 512 bytes
      isLikelyText(source)          -> detectTextSubtype(source)
      else                          -> FileFormat.BINARY
    }
  }

  private fun isLikelyText(source: MmapSource): Boolean {
    val sample = source.readBytes(0, minOf(512, source.size.toInt()))
    val nonPrintable = sample.count { b ->
      val i = b.toInt() and 0xFF
      i < 9 || (i in 14..31 && i != 27)  // not tab/LF/CR
    }
    return nonPrintable.toDouble() / sample.size < 0.02  // <2% non-printable
  }
}
\end{lstlisting}

\PB

% ============================================================
% CHAPTER 3 — PDF ENGINE
% ============================================================
\chapter{PDF Engine — Sub-80ms First Paint}

\section{Why Android's PdfRenderer Is Enough}

\texttt{android.graphics.pdf.PdfRenderer} (API 21+) is a wrapper around the
same PDFium engine used by Google Chrome. It is hardware-accelerated, supports
the full PDF 1.7 specification including forms, encryption (with password), and
embedded fonts. It requires \textbf{zero} external libraries.

The reason nobody uses it for fast viewing is that they call:
\begin{lstlisting}[language=Kotlin]
page.render(bitmap, null, null, PdfRenderer.Page.RENDER_MODE_FOR_DISPLAY)
\end{lstlisting}
where \texttt{bitmap} is the full page at 150 DPI — a 2480×3508 pixel bitmap
= 35MB RAM per page, render time 800ms--2s. This is the mistake we avoid.

\section{FLUX PDF Strategy: Tile-Based Progressive Rendering}

\begin{lstlisting}[language=Kotlin,caption={FluxPdfRenderer — tile-based PDF viewer}]
class FluxPdfRenderer(private val file: File) : AutoCloseable {
  private val fd       = ParcelFileDescriptor.open(file,
                           ParcelFileDescriptor.MODE_READ_ONLY)
  private val renderer = PdfRenderer(fd)
  private val pageCount get() = renderer.pageCount

  // Page size at 72 DPI (PDF native resolution)
  fun getPageSize(pageIndex: Int): SizeF {
    val page = renderer.openPage(pageIndex)
    val size = SizeF(page.width.toFloat(), page.height.toFloat())
    page.close()
    return size
  }

  // Render ONLY the visible tile region of a page
  // clipRect is in PDF point coordinates (72 DPI)
  fun renderRegion(
    pageIndex:   Int,
    clipRect:    Rect,   // which region to render (PDF coords)
    outputW:     Int,    // output bitmap width (screen pixels)
    outputH:     Int,    // output bitmap height (screen pixels)
  ): Bitmap {
    val bitmap = Bitmap.createBitmap(outputW, outputH, Bitmap.Config.RGB_565)
    // RGB_565 not ARGB_8888: 2 bytes/px not 4. Same quality for PDFs.

    renderer.openPage(pageIndex).use { page ->
      // The transform maps clipRect in PDF coords to the output bitmap
      val scaleX = outputW.toFloat() / clipRect.width()
      val scaleY = outputH.toFloat() / clipRect.height()
      val transform = Matrix().apply {
        setScale(scaleX, scaleY)
        postTranslate(-clipRect.left * scaleX, -clipRect.top * scaleY)
      }
      page.render(bitmap, null, transform,
        PdfRenderer.Page.RENDER_MODE_FOR_DISPLAY)
    }
    return bitmap
  }

  override fun close() { renderer.close(); fd.close() }
}

// TileRenderer implementation for PDF
class PdfTileRenderer(
  private val pdf:       FluxPdfRenderer,
  private val pageIndex: Int,
  private val pageSize:  SizeF,
  private val screenDPI: Float = 160f,
) : TileRenderer {

  override fun renderTile(key: TileKey, tileW: Int, tileH: Int): Bitmap {
    // Scale factor: how many screen pixels per PDF point
    val scale = screenDPI / 72f * when(key.zoomLevel) {
      0 -> 0.5f   // thumbnail zoom
      1 -> 1.0f   // fit-width zoom
      2 -> 2.0f   // 2x zoom
      else -> 1.0f
    }
    // Convert tile position back to PDF coordinate space
    val pdfLeft   = (key.col * tileW / scale).toInt()
    val pdfTop    = (key.row * tileH / scale).toInt()
    val pdfRight  = minOf(((key.col + 1) * tileW / scale).toInt(), pageSize.width.toInt())
    val pdfBottom = minOf(((key.row + 1) * tileH / scale).toInt(), pageSize.height.toInt())

    return pdf.renderRegion(
      pageIndex = pageIndex,
      clipRect  = Rect(pdfLeft, pdfTop, pdfRight, pdfBottom),
      outputW   = tileW,
      outputH   = tileH,
    )
  }
}
\end{lstlisting}

\section{PDF Page Prefetching Strategy}

\begin{lstlisting}[language=Kotlin,caption={Smart prefetch: render adjacent pages at low res}]
class PdfPrefetcher(private val pdf: FluxPdfRenderer) {
  // When user is on page N, prefetch N-1 and N+1 at thumbnail quality
  // This makes page-turn feel instant
  fun prefetchAdjacent(currentPage: Int, pageCount: Int) {
    val toFetch = mutableListOf<Int>()
    if (currentPage > 0)            toFetch += currentPage - 1
    if (currentPage < pageCount-1)  toFetch += currentPage + 1

    toFetch.forEach { pageIdx ->
      if (thumbnailCache[pageIdx] == null) {
        backgroundPool.submit {
          val size    = pdf.getPageSize(pageIdx)
          val thumbW  = 360  // 360px width thumbnail
          val thumbH  = (360 * size.height / size.width).toInt()
          val thumb   = pdf.renderRegion(
            pageIndex = pageIdx,
            clipRect  = Rect(0, 0, size.width.toInt(), size.height.toInt()),
            outputW   = thumbW,
            outputH   = thumbH,
          )
          thumbnailCache[pageIdx] = thumb
          onThumbnailReady(pageIdx, thumb)
        }
      }
    }
  }
}
\end{lstlisting}

\section{PDF Performance Targets}

\begin{center}\ra{1.35}
\begin{tabular}{L{4.8cm} C{2.4cm} C{2.4cm} L{4.2cm}}
\TH{Operation} & \TH{Current Apps} & \TH{FLUX Target} & \TH{How We Achieve It} \\
\rowcolor{rowA} First page visible & 1.5--3\,s & \textbf{$<$80ms} & Render first tile at 72 DPI immediately \\
\rowcolor{rowB} Full first page at read DPI & 800ms--2\,s & \textbf{$<$200ms} & Tile-based: visible region only \\
\rowcolor{rowA} Page turn (adjacent) & 300ms--1\,s & \textbf{$<$50ms} & Pre-fetched thumbnail already cached \\
\rowcolor{rowB} Search text in 100-page PDF & 5--30\,s & \textbf{$<$500ms} & \texttt{PdfRenderer.Page.openPage().getTextContents()} streaming \\
\rowcolor{rowA} RAM at 10-page PDF & 50--350\,MB & \textbf{$<$8\,MB} & Only 2--3 tiles per page in cache at once \\
\rowcolor{rowB} RAM at 100-page PDF & 200MB--1GB & \textbf{$<$15\,MB} & Same tile cache, page count doesn't matter \\
\end{tabular}
\end{center}

\section{PDF Editing — Annotation Layer}

\begin{lstlisting}[language=Kotlin,caption={PDF annotation overlay (no library needed)}]
// Annotations are stored as an in-memory overlay
// Written to PDF using PdfDocument API on save
data class PdfAnnotation(
  val pageIndex:  Int,
  val type:       AnnotationType,  // HIGHLIGHT, UNDERLINE, TEXT, DRAW
  val rect:       RectF,           // in PDF point coordinates
  val color:      Int,
  val content:    String = "",     // for TEXT annotations
  val path:       List<PointF> = emptyList() // for DRAW annotations
)

enum class AnnotationType { HIGHLIGHT, UNDERLINE, STRIKETHROUGH, TEXT, DRAW, STAMP }

class PdfAnnotationLayer {
  private val annotations = mutableListOf<PdfAnnotation>()

  fun addHighlight(pageIndex: Int, rect: RectF, color: Int = 0xFFFFFF00.toInt()) {
    annotations += PdfAnnotation(pageIndex, AnnotationType.HIGHLIGHT, rect, color)
  }

  // Render annotations on top of the tile bitmap
  fun drawAnnotations(pageIndex: Int, canvas: Canvas, pdfToScreenMatrix: Matrix) {
    val paint = Paint().apply { isAntiAlias = true; style = Paint.Style.FILL }
    annotations.filter { it.pageIndex == pageIndex }.forEach { ann ->
      val screenRect = RectF(ann.rect).also { pdfToScreenMatrix.mapRect(it) }
      when (ann.type) {
        AnnotationType.HIGHLIGHT -> {
          paint.color = (ann.color and 0x00FFFFFF) or 0x60000000 // 38% alpha
          canvas.drawRect(screenRect, paint)
        }
        AnnotationType.DRAW -> {
          paint.style = Paint.Style.STROKE
          paint.strokeWidth = 3f
          paint.color = ann.color
          val path = Path()
          ann.path.forEachIndexed { i, pt ->
            val screenPt = floatArrayOf(pt.x, pt.y)
            pdfToScreenMatrix.mapPoints(screenPt)
            if (i == 0) path.moveTo(screenPt[0], screenPt[1])
            else        path.lineTo(screenPt[0], screenPt[1])
          }
          canvas.drawPath(path, paint)
        }
        else -> { /* other annotation types */ }
      }
    }
  }

  // Save annotations into new PDF using PdfDocument
  suspend fun saveToPdf(originalFile: File, outputFile: File) =
    withContext(Dispatchers.IO) {
      // Read original pages, re-render + overlay annotations
      // Write to output using android.graphics.pdf.PdfDocument
      PdfDocument().use { doc ->
        PdfRenderer(ParcelFileDescriptor.open(originalFile,
          ParcelFileDescriptor.MODE_READ_ONLY)).use { renderer ->
          for (i in 0 until renderer.pageCount) {
            renderer.openPage(i).use { page ->
              val pageInfo = PdfDocument.PageInfo.Builder(
                page.width, page.height, i + 1).create()
              val docPage  = doc.startPage(pageInfo)
              page.render(docPage.canvas, null, null,
                PdfRenderer.Page.RENDER_MODE_FOR_PRINT)
              // Draw annotations on top
              drawAnnotations(i, docPage.canvas, Matrix()) // identity matrix for save
              doc.finishPage(docPage)
            }
          }
          outputFile.outputStream().use { doc.writeTo(it) }
        }
      }
    }
}
\end{lstlisting}

\PB

% ============================================================
% CHAPTER 4 — DOCX ENGINE
% ============================================================
\chapter{DOCX/DOC Engine — Minimal OOXML Parser}

\section{What DOCX Actually Is}

A DOCX file is a ZIP archive. Unzip any DOCX and you find:

\begin{center}\ra{1.3}
\begin{tabular}{L{4.0cm} L{9.8cm}}
\TH{File Inside ZIP} & \TH{Contents} \\
\rowcolor{rowA} \texttt{word/document.xml} & The main document body: paragraphs, runs, tables \\
\rowcolor{rowB} \texttt{word/styles.xml} & Style definitions: Heading1, Normal, etc. \\
\rowcolor{rowA} \texttt{word/fontTable.xml} & Font references \\
\rowcolor{rowB} \texttt{word/numbering.xml} & List/bullet definitions \\
\rowcolor{rowA} \texttt{word/media/image1.jpg} & Embedded images \\
\rowcolor{rowB} \texttt{word/\_rels/document.xml.rels} & Relationship map \\
\rowcolor{rowA} \texttt{[Content\_Types].xml} & MIME type map for each part \\
\end{tabular}
\end{center}

\texttt{document.xml} uses these XML elements for 95\% of real-world documents:
\texttt{w:p} (paragraph), \texttt{w:r} (run/text span), \texttt{w:t} (text
content), \texttt{w:tbl} (table), \texttt{w:tr} (table row), \texttt{w:tc}
(table cell), \texttt{w:pPr} (paragraph properties), \texttt{w:rPr} (run
properties: bold, italic, size, colour), \texttt{w:drawing} (image).

\section{FLUX DocxParser — Zero Library}

\begin{lstlisting}[language=Kotlin,caption={Streaming DOCX parser using built-in XML}]
// OOXML namespace prefix
const val W = "http://schemas.openxmlformats.org/wordprocessingml/2006/main"

data class DocParagraph(
  val runs:      List<DocRun>,
  val styleId:   String,    // "Heading1", "Normal", etc.
  val alignment: Int,       // Gravity.LEFT / CENTER / RIGHT / END
  val spaceBefore: Int,     // twips
  val spaceAfter:  Int,
  val listLevel: Int = -1,  // -1 = not a list item
)

data class DocRun(
  val text:      String,
  val bold:      Boolean = false,
  val italic:    Boolean = false,
  val underline: Boolean = false,
  val strikethrough: Boolean = false,
  val fontSize:  Int = 24,    // half-points (24 = 12pt)
  val color:     Int = 0xFF000000.toInt(),
  val highlight: Int? = null,
)

sealed class DocElement {
  data class Paragraph(val para: DocParagraph) : DocElement()
  data class Table(val rows: List<List<DocParagraph>>) : DocElement()
  data class Image(val rid: String, val width: Int, val height: Int) : DocElement()
  data class PageBreak(val dummy: Unit = Unit) : DocElement()
}

class DocxParser {
  // Parse DOCX from ZipInputStream — streaming, no full-file load
  fun parse(file: File): DocxDocument {
    val elements  = mutableListOf<DocElement>()
    val images    = mutableMapOf<String, ByteArray>()  // rId -> bytes
    val styles    = mutableMapOf<String, DocStyle>()

    ZipFile(file).use { zip ->
      // 1. Parse styles first
      zip.getEntry("word/styles.xml")?.let { entry ->
        parseStyles(zip.getInputStream(entry), styles)
      }
      // 2. Parse relationships (for image rIds)
      val rels = mutableMapOf<String, String>()
      zip.getEntry("word/_rels/document.xml.rels")?.let { entry ->
        parseRels(zip.getInputStream(entry), rels)
      }
      // 3. Stream-parse document.xml — the main content
      zip.getEntry("word/document.xml")?.let { entry ->
        streamParseDocument(zip.getInputStream(entry), styles, elements)
      }
      // 4. Load embedded images lazily (only those referenced)
      rels.forEach { (rId, target) ->
        if (target.contains("media/")) {
          zip.getEntry("word/$target")?.let { entry ->
            images[rId] = zip.getInputStream(entry).readBytes()
          }
        }
      }
    }
    return DocxDocument(elements, images)
  }

  private fun streamParseDocument(
    input:    InputStream,
    styles:   Map<String, DocStyle>,
    elements: MutableList<DocElement>
  ) {
    val factory = XmlPullParserFactory.newInstance()
    val parser  = factory.newPullParser()
    parser.setInput(input, "UTF-8")

    var currentPara     : MutableList<DocRun>? = null
    var currentRun      : StringBuilder? = null
    var currentRunProps  = DocRunProps()
    var currentParaProps = DocParaProps()
    var inTable          = false
    val tableRows        = mutableListOf<MutableList<DocParagraph>>()
    val currentRow       = mutableListOf<DocParagraph>()

    var eventType = parser.next()
    while (eventType != XmlPullParser.END_DOCUMENT) {
      when (eventType) {
        XmlPullParser.START_TAG -> when (parser.name) {
          "w:p"   -> { currentPara = mutableListOf(); currentParaProps = DocParaProps() }
          "w:r"   -> { currentRun = StringBuilder(); currentRunProps = DocRunProps() }
          "w:t"   -> { currentRun?.append(parser.nextText()) }
          "w:b"   -> currentRunProps = currentRunProps.copy(bold = true)
          "w:i"   -> currentRunProps = currentRunProps.copy(italic = true)
          "w:u"   -> currentRunProps = currentRunProps.copy(underline = true)
          "w:sz"  -> currentRunProps = currentRunProps.copy(
            fontSize = parser.getAttributeValue(W, "val")?.toIntOrNull() ?: 24)
          "w:color" -> {
            val hex = parser.getAttributeValue(W, "val") ?: "000000"
            if (hex != "auto") currentRunProps = currentRunProps.copy(
              color = (0xFF000000.toInt() or hex.toLong(16).toInt()))
          }
          "w:jc"  -> {
            val align = parser.getAttributeValue(W, "val") ?: "left"
            currentParaProps = currentParaProps.copy(alignment = when(align) {
              "center" -> Gravity.CENTER_HORIZONTAL
              "right"  -> Gravity.END
              "both"   -> Gravity.FILL_HORIZONTAL
              else     -> Gravity.START
            })
          }
          "w:pStyle" -> {
            val sId = parser.getAttributeValue(W, "val") ?: "Normal"
            currentParaProps = currentParaProps.copy(styleId = sId)
          }
          "w:tbl" -> { inTable = true; tableRows.clear() }
          "w:tr"  -> currentRow.clear()
        }
        XmlPullParser.END_TAG -> when (parser.name) {
          "w:r" -> {
            currentRun?.toString()?.takeIf { it.isNotEmpty() }?.let { text ->
              currentPara?.add(DocRun(
                text      = text,
                bold      = currentRunProps.bold,
                italic    = currentRunProps.italic,
                underline = currentRunProps.underline,
                fontSize  = currentRunProps.fontSize,
                color     = currentRunProps.color,
              ))
            }
            currentRun = null
          }
          "w:p" -> {
            currentPara?.let { runs ->
              val para = DocParagraph(
                runs      = runs,
                styleId   = currentParaProps.styleId,
                alignment = currentParaProps.alignment,
                spaceBefore = currentParaProps.spaceBefore,
                spaceAfter  = currentParaProps.spaceAfter,
              )
              if (inTable) currentRow.add(para)
              else         elements.add(DocElement.Paragraph(para))
            }
            currentPara = null
          }
          "w:tbl" -> {
            elements.add(DocElement.Table(tableRows.map { it.toList() }))
            inTable = false
          }
        }
      }
      eventType = parser.next()
    }
  }
}
\end{lstlisting}

\section{DOCX Layout Engine — Canvas-Based Rendering}

\begin{lstlisting}[language=Kotlin,caption={DocxLayoutEngine — text layout without TextView}]
class DocxLayoutEngine(
  private val doc:       DocxDocument,
  private val pageWidth: Int,   // pixels
  private val margin:    Int = 48, // pixels
) {
  private val textWidth = pageWidth - 2 * margin
  private val linePaint = Paint().apply { isAntiAlias = true }

  data class LayoutLine(
    val y:      Float,   // baseline y position in the document canvas
    val spans:  List<LayoutSpan>,
  )
  data class LayoutSpan(val x: Float, val text: String, val paint: Paint)

  fun layoutAll(): List<LayoutLine> {
    val lines = mutableListOf<LayoutLine>()
    var y     = margin.toFloat()

    doc.elements.forEach { element ->
      when (element) {
        is DocElement.Paragraph -> {
          y = layoutParagraph(element.para, lines, y)
          y += element.para.spaceAfter / 20f  // twips to px (approx)
        }
        is DocElement.Table -> {
          y = layoutTable(element, lines, y)
        }
        is DocElement.PageBreak -> {
          y = nextPageY(y)
        }
        else -> {}
      }
    }
    return lines
  }

  private fun layoutParagraph(
    para:  DocParagraph,
    lines: MutableList<LayoutLine>,
    startY: Float
  ): Float {
    var y      = startY
    val words  = mutableListOf<Pair<String, Paint>>() // word + its paint

    // Flatten runs into word list
    para.runs.forEach { run ->
      val paint = Paint().apply {
        isAntiAlias = true
        typeface    = resolveFace(run.bold, run.italic)
        textSize    = run.fontSize / 2f * (displayDPI / 72f)  // half-pt to px
        color       = run.color
      }
      run.text.split(" ").forEach { word ->
        if (word.isNotEmpty()) words += Pair("$word ", paint)
      }
    }

    // Greedy line-breaking
    var lineStart   = 0
    var lineWidth   = 0f
    var currentLine = mutableListOf<LayoutSpan>()
    var x           = margin.toFloat()

    words.forEachIndexed { i, (word, paint) ->
      val wordW = paint.measureText(word)
      if (x + wordW > margin + textWidth && currentLine.isNotEmpty()) {
        // Flush current line
        lines.add(LayoutLine(y + paint.fontMetrics.ascent.unaryMinus(), currentLine))
        y          += paint.fontMetrics.let { it.descent - it.ascent + it.leading }
        currentLine = mutableListOf()
        x           = margin.toFloat()
      }
      currentLine.add(LayoutSpan(x, word, paint))
      x += wordW
    }
    if (currentLine.isNotEmpty()) {
      val paint = currentLine.last().paint
      lines.add(LayoutLine(y + paint.fontMetrics.ascent.unaryMinus(), currentLine))
      y += paint.fontMetrics.let { it.descent - it.ascent + it.leading }
    }
    return y
  }
}
\end{lstlisting}

\section{DOCX Performance Targets}

\begin{center}\ra{1.35}
\begin{tabular}{L{4.4cm} C{2.6cm} C{2.4cm} L{4.4cm}}
\TH{Operation} & \TH{WPS Office} & \TH{FLUX Target} & \TH{How} \\
\rowcolor{rowA} First page visible & 3--6\,s & \textbf{$<$150ms} & Parse first 50 elements, layout first viewport only \\
\rowcolor{rowB} Full 20-page parse & 3--6\,s & \textbf{$<$400ms} & SAX streaming parser, no DOM tree \\
\rowcolor{rowA} RAM for 10-page doc & 80--200\,MB & \textbf{$<$12\,MB} & Only visible lines in canvas buffer \\
\rowcolor{rowB} Edit single word & 500ms re-render & \textbf{$<$16ms} & Re-layout only the affected paragraph \\
\rowcolor{rowA} Save edited DOCX & 2--5\,s & \textbf{$<$300ms} & Patch only modified XML nodes in ZIP \\
\end{tabular}
\end{center}

\PB

% ============================================================
% CHAPTER 5 — XLSX ENGINE
% ============================================================
\chapter{XLSX Engine — Virtual Grid Renderer}

\section{What XLSX Actually Is}

An XLSX file is also a ZIP containing XML. The key files:

\begin{center}\ra{1.3}
\begin{tabular}{L{4.2cm} L{9.6cm}}
\TH{File} & \TH{Contents} \\
\rowcolor{rowA} \texttt{xl/workbook.xml} & Sheet names and references \\
\rowcolor{rowB} \texttt{xl/worksheets/sheet1.xml} & Cell data: addresses, values, formula references \\
\rowcolor{rowA} \texttt{xl/sharedStrings.xml} & String lookup table (all string values stored once here) \\
\rowcolor{rowB} \texttt{xl/styles.xml} & Cell format definitions (number format, font, fill, border) \\
\rowcolor{rowA} \texttt{xl/calcChain.xml} & Formula calculation order \\
\end{tabular}
\end{center}

\section{Virtual Grid — Key Insight}

A spreadsheet with 10,000 rows × 50 columns has 500,000 cells. Rendering all
of them at once means 500,000 \texttt{drawText()} calls per frame — 50 FPS
impossible. The solution: \textbf{virtual grid rendering}. Draw only the cells
visible in the current viewport, exactly like \texttt{RecyclerView} but for a
2D grid drawn directly on Canvas.

\begin{lstlisting}[language=Kotlin,caption={XlsxVirtualGrid — O(visible cells) per frame}]
class XlsxVirtualGrid(private val sheet: SheetData) {
  // Column widths and row heights (from xlsx styles, or default)
  private val colWidths = FloatArray(sheet.maxCol + 1) { DEFAULT_COL_W }
  private val rowHeights= FloatArray(sheet.maxRow + 1) { DEFAULT_ROW_H }

  // Cumulative positions for O(log n) binary search
  private val colX = FloatArray(sheet.maxCol + 2)  // colX[i] = pixel x of column i
  private val rowY = FloatArray(sheet.maxRow + 2)

  init {
    // Build cumulative arrays once — O(n) total, O(log n) per viewport query
    colX[0] = HEADER_W  // row header width
    for (c in 0..sheet.maxCol) colX[c+1] = colX[c] + colWidths[c]
    rowY[0] = HEADER_H  // column header height
    for (r in 0..sheet.maxRow) rowY[r+1] = rowY[r] + rowHeights[r]
  }

  // Find first visible column — O(log n)
  private fun firstVisibleCol(scrollX: Float): Int {
    var lo = 0; var hi = sheet.maxCol
    while (lo < hi) {
      val mid = (lo + hi) / 2
      if (colX[mid + 1] < scrollX) lo = mid + 1 else hi = mid
    }
    return lo
  }

  fun onDraw(canvas: Canvas, viewport: RectF, scrollX: Float, scrollY: Float) {
    val firstCol = firstVisibleCol(scrollX)     // O(log n)
    val firstRow = firstVisibleRow(scrollY)     // O(log n)
    val lastCol  = lastVisibleCol(scrollX + viewport.width(), firstCol)
    val lastRow  = lastVisibleRow(scrollY + viewport.height(), firstRow)

    val cellPaint = Paint().apply { isAntiAlias = true }
    val textPaint = Paint().apply {
      isAntiAlias = true; textSize = 13f * displayScale
    }
    val gridPaint = Paint().apply { color = 0xFFD0D0D0.toInt(); strokeWidth = 1f }

    // Draw only visible cells — O(visible_rows × visible_cols)
    for (r in firstRow..lastRow) {
      for (c in firstCol..lastCol) {
        val x   = colX[c] - scrollX + viewport.left
        val y   = rowY[r] - scrollY + viewport.top
        val w   = colWidths[c]
        val h   = rowHeights[r]
        val cell = sheet.getCell(r, c)  // O(1) HashMap lookup

        // Cell background
        val fillColor = cell?.style?.fillColor ?: 0xFFFFFFFF.toInt()
        if (fillColor != 0xFFFFFFFF.toInt()) {
          cellPaint.color = fillColor
          canvas.drawRect(x, y, x+w, y+h, cellPaint)
        }

        // Cell text
        cell?.displayValue?.let { text ->
          textPaint.color    = cell.style?.textColor ?: 0xFF000000.toInt()
          textPaint.isFakeBoldText = cell.style?.bold ?: false
          textPaint.textAlign = when(cell.style?.align) {
            "right"  -> Paint.Align.RIGHT
            "center" -> Paint.Align.CENTER
            else     -> Paint.Align.LEFT
          }
          // Clip to cell bounds to prevent overflow
          canvas.save()
          canvas.clipRect(x, y, x+w, y+h)
          canvas.drawText(text, x + 4, y + h - 5, textPaint)
          canvas.restore()
        }

        // Grid lines
        canvas.drawLine(x+w, y, x+w, y+h, gridPaint) // vertical
        canvas.drawLine(x, y+h, x+w, y+h, gridPaint) // horizontal
      }
    }

    // Frozen row/column headers
    drawHeaders(canvas, firstCol, lastCol, firstRow, lastRow, scrollX, scrollY)
  }
}
\end{lstlisting}

\section{XLSX Cell Parser — Shared Strings Table}

\begin{lstlisting}[language=Kotlin,caption={SheetData — parse cells with shared strings}]
class SheetData {
  val cells = HashMap<Long, CellData>()  // key = row * MAX_COL + col

  // O(1) cell lookup
  fun getCell(row: Int, col: Int): CellData? =
    cells[row.toLong() * MAX_COL + col]

  companion object {
    fun parse(xlsxFile: File, sheetIndex: Int = 0): SheetData {
      val data = SheetData()
      ZipFile(xlsxFile).use { zip ->
        // 1. Load shared strings table into ArrayList
        val sharedStrings = ArrayList<String>(1000)
        zip.getEntry("xl/sharedStrings.xml")?.let { entry ->
          val parser = XmlPullParserFactory.newInstance().newPullParser()
          parser.setInput(zip.getInputStream(entry), "UTF-8")
          var inT = false
          var event = parser.next()
          while (event != XmlPullParser.END_DOCUMENT) {
            when {
              event == XmlPullParser.START_TAG && parser.name == "t" -> {
                sharedStrings.add(parser.nextText())
              }
            }
            event = parser.next()
          }
        }

        // 2. Stream-parse sheet XML
        zip.getEntry("xl/worksheets/sheet${sheetIndex+1}.xml")?.let { entry ->
          val parser = XmlPullParserFactory.newInstance().newPullParser()
          parser.setInput(zip.getInputStream(entry), "UTF-8")
          var currentRef  = ""
          var currentType = ""  // "s"=sharedString, "n"=number, "b"=boolean
          var inV = false
          var event = parser.next()
          while (event != XmlPullParser.END_DOCUMENT) {
            when (event) {
              XmlPullParser.START_TAG -> when (parser.name) {
                "c" -> {
                  currentRef  = parser.getAttributeValue(null, "r") ?: ""
                  currentType = parser.getAttributeValue(null, "t") ?: "n"
                }
                "v" -> inV = true
              }
              XmlPullParser.TEXT -> if (inV) {
                val rawValue = parser.text
                val (row, col) = parseCellRef(currentRef)
                val displayValue = when (currentType) {
                  "s"  -> sharedStrings.getOrElse(rawValue.toInt()) { rawValue }
                  "b"  -> if (rawValue == "1") "TRUE" else "FALSE"
                  else -> rawValue  // number: display as-is or format
                }
                data.cells[row.toLong() * MAX_COL + col] =
                  CellData(row, col, rawValue, displayValue)
              }
              XmlPullParser.END_TAG -> if (parser.name == "v") inV = false
            }
            event = parser.next()
          }
        }
      }
      return data
    }

    // "C5" -> (row=4, col=2)
    fun parseCellRef(ref: String): Pair<Int, Int> {
      val colStr = ref.takeWhile { it.isLetter() }
      val rowStr = ref.dropWhile { it.isLetter() }
      var col = 0
      colStr.uppercase().forEach { c -> col = col * 26 + (c - 'A' + 1) }
      return Pair(rowStr.toInt() - 1, col - 1)
    }
  }
}
\end{lstlisting}

\PB

% ============================================================
% CHAPTER 6 — IMAGE ENGINE
% ============================================================
\chapter{Image Engine — BitmapRegionDecoder + Progressive Tiling}

\section{The Core API: BitmapRegionDecoder}

\texttt{android.graphics.BitmapRegionDecoder} (API 10+) decodes a
\emph{subregion} of a JPEG, PNG, or WebP image. It works by seeking to the
correct compressed data offset and decoding only the necessary MCU blocks
(for JPEG) or IDAT chunks (for PNG). The operating system's image codec
handles this at native C++ speed.

\begin{lstlisting}[language=Kotlin,caption={ImageTileRenderer using BitmapRegionDecoder}]
class ImageTileRenderer(file: File) : TileRenderer, AutoCloseable {
  private val decoder: BitmapRegionDecoder =
    BitmapRegionDecoder.newInstance(file.absolutePath, false)
      ?: error("Cannot decode: ${file.name}")

  val imageWidth:  Int = decoder.width
  val imageHeight: Int = decoder.height

  private val options = BitmapFactory.Options().apply {
    inPreferredConfig = Bitmap.Config.RGB_565  // 50% less RAM than ARGB_8888
    inSampleSize      = 1
  }

  override fun renderTile(key: TileKey, tileW: Int, tileH: Int): Bitmap {
    // Map tile position to image pixel region
    val scale     = zoomLevelToScale(key.zoomLevel)
    val imgLeft   = (key.col * tileW / scale).toInt().coerceIn(0, imageWidth)
    val imgTop    = (key.row * tileH / scale).toInt().coerceIn(0, imageHeight)
    val imgRight  = ((key.col + 1) * tileW / scale).toInt().coerceIn(0, imageWidth)
    val imgBottom = ((key.row + 1) * tileH / scale).toInt().coerceIn(0, imageHeight)

    if (imgLeft >= imgRight || imgTop >= imgBottom) {
      return Bitmap.createBitmap(tileW, tileH, Bitmap.Config.RGB_565)
    }

    // BitmapRegionDecoder inSampleSize: use nearest power-of-2 downscale
    options.inSampleSize = nearestPow2DownSample(
      srcW = imgRight - imgLeft,  srcH = imgBottom - imgTop,
      dstW = tileW,               dstH = tileH
    )

    return decoder.decodeRegion(
      Rect(imgLeft, imgTop, imgRight, imgBottom), options
    )
    // Only this rectangle is decoded from the compressed file.
    // For a 100MP image (10000x10000px): decoding one 512x512 tile
    // takes ~5-15ms and uses ~0.5MB RAM — not 400MB.
  }

  private fun nearestPow2DownSample(srcW:Int,srcH:Int,dstW:Int,dstH:Int): Int {
    var sample = 1
    while (srcW / (sample * 2) >= dstW && srcH / (sample * 2) >= dstH) {
      sample *= 2
    }
    return sample
  }

  override fun close() = decoder.recycle()
}
\end{lstlisting}

\section{Image Format Support Matrix}

\begin{center}\ra{1.35}
\begin{tabular}{L{2.4cm} C{1.8cm} C{2.0cm} C{2.4cm} L{5.2cm}}
\TH{Format} & \TH{BitmapRegionDecoder} & \TH{Android API} & \TH{RAM for 12MP} & \TH{Strategy} \\
\rowcolor{rowA} JPEG & \textbf{Yes} & 10+ & $\sim$0.5MB/tile & BitmapRegionDecoder, MCU-aligned tiles \\
\rowcolor{rowB} PNG & \textbf{Yes} & 10+ & $\sim$0.5MB/tile & BitmapRegionDecoder \\
\rowcolor{rowA} WebP & \textbf{Yes} & 14+ & $\sim$0.5MB/tile & BitmapRegionDecoder \\
\rowcolor{rowB} HEIC/HEIF & No & 28+ & Full decode needed & \texttt{ImageDecoder} API (API 28+) with sample size \\
\rowcolor{rowA} GIF & No & — & Varies & \texttt{ImageDecoder} + \texttt{AnimatedImageDrawable} \\
\rowcolor{rowB} RAW (DNG) & No & — & 50--100MB & \texttt{DngCreator} / external thumbnail first; full decode async \\
\rowcolor{rowA} BMP & No & — & Full decode & \texttt{BitmapFactory} with \texttt{inSampleSize} \\
\rowcolor{rowB} SVG & No & — & Vector, tiny & Custom \texttt{Canvas} parser (Chapter 11) \\
\end{tabular}
\end{center}

\section{Progressive Loading Strategy}

\begin{center}\ra{1.3}
\begin{tabular}{C{0.8cm} L{2.4cm} L{2.4cm} C{2.0cm} L{5.8cm}}
\TH{Stage} & \TH{What Rendered} & \TH{When} & \TH{Latency} & \TH{How} \\
\rowcolor{rowA} 1 & Thumbnail (96×96) & Immediately & 0ms & From FLUX ThumbnailCache \\
\rowcolor{rowB} 2 & Low-res full image & 30--80ms & 50--80ms & BitmapRegionDecoder, inSampleSize=8 \\
\rowcolor{rowA} 3 & Full-res visible region & After stage 2 & 10--30ms & BitmapRegionDecoder, inSampleSize=1, clip to viewport \\
\rowcolor{rowB} 4 & Full-res with zoom & On user pinch & 5--20ms/tile & TileEngine fetches new tiles at new zoom level \\
\end{tabular}
\end{center}

\PB

% ============================================================
% CHAPTER 7 — TEXT / CODE ENGINE
% ============================================================
\chapter{Text \& Code Engine — Canvas-Based Fast Renderer}

\section{Why Not TextView}

Android's \texttt{TextView} loads the \emph{entire text} into a
\texttt{StaticLayout} before displaying anything. For a 10\,MB source file
(approximately 300,000 lines of code), \texttt{StaticLayout} takes 3--8
seconds to compute layout and uses 80--200\,MB RAM. The view then
becomes unresponsive during scrolling.

Our solution: \textbf{draw text directly on Canvas, using a line-virtual
approach}. We compute layout only for visible lines, exactly like
\texttt{RecyclerView} but at the character level.

\section{FluxTextRenderer}

\begin{lstlisting}[language=Kotlin,caption={FluxTextRenderer — O(visible lines) per frame}]
class FluxTextRenderer(
  file:        File,
  private val textSize:    Float = 14f,
  private val lineSpacing: Float = 1.4f,
) {
  // Memory-mapped source — never loads full file
  private val source     = MmapSource(file)
  private val lineIndex  = LineIndex(source)   // built lazily
  private val totalLines get() = lineIndex.lineCount

  // Line height in pixels
  private val lineH: Float
  private val textPaint = Paint().apply {
    isAntiAlias = true
    typeface    = Typeface.MONOSPACE
    this.textSize = textSize * displayScale
    color       = 0xFF212121.toInt()
  }
  init { lineH = textPaint.fontMetrics.let { (it.descent - it.ascent) * lineSpacing } }

  fun onDraw(canvas: Canvas, scrollY: Float, viewportH: Float) {
    val firstLine = (scrollY / lineH).toInt().coerceAtLeast(0)
    val lastLine  = ((scrollY + viewportH) / lineH).toInt()
      .coerceAtMost(totalLines - 1)

    for (lineNum in firstLine..lastLine) {
      val y     = lineNum * lineH - scrollY + lineH  // baseline y
      val line  = lineIndex.getLine(lineNum)          // O(1): mmap offset lookup

      // Syntax highlighting: tokenize and draw spans
      val spans = highlighter.highlight(line)
      var x     = GUTTER_W + PADDING

      spans.forEach { span ->
        textPaint.color = span.color
        canvas.drawText(span.text, x, y, textPaint)
        x += textPaint.measureText(span.text)
      }
    }

    // Line number gutter
    drawGutter(canvas, firstLine, lastLine, scrollY)
  }

  val contentHeight: Float get() = totalLines * lineH
  val contentWidth:  Float get() = lineIndex.maxLineWidth * charWidth
}

// Builds an index of line start offsets in the mmap — never loads lines
class LineIndex(private val source: MmapSource) {
  // Compact array of line-start byte offsets
  // For 300k lines × 4 bytes = 1.2MB — acceptable
  private val offsets = IntArray(estimateLineCount(source) + 1)
  private var count   = 0

  init { buildIndex() }

  private fun buildIndex() {
    offsets[count++] = 0  // line 0 starts at byte 0
    for (i in 0 until source.size) {
      if (source.readByte(i) == '\n'.code.toByte()) {
        offsets[count++] = (i + 1).toInt()
      }
    }
  }

  val lineCount: Int get() = count

  // Get line as String — reads only that line's bytes from mmap
  fun getLine(lineNum: Int): String {
    if (lineNum >= count) return ""
    val start  = offsets[lineNum].toLong()
    val end    = if (lineNum + 1 < count) offsets[lineNum + 1].toLong()
                 else source.size
    val length = (end - start).toInt().coerceAtLeast(0)
    if (length == 0) return ""
    return String(source.readBytes(start, minOf(length, MAX_LINE_LEN)), Charsets.UTF_8)
      .trimEnd('\r', '\n')
  }

  var maxLineWidth = 0
  private fun estimateLineCount(s: MmapSource): Int = (s.size / 40).toInt()
}
\end{lstlisting}

\section{Syntax Highlighter — Regex Token-Based}

\begin{lstlisting}[language=Kotlin,caption={Syntax highlighter — no AST, pure regex spans}]
class SyntaxHighlighter(private val language: Language) {
  data class TextSpan(val text: String, val color: Int)

  // Language-specific token patterns — ordered by priority
  private val patterns: List<Pair<Regex, Int>> = when (language) {
    Language.KOTLIN, Language.JAVA -> listOf(
      Regex("""//.*$""") to 0xFF78909C.toInt(),           // comment
      Regex(""""([^"\\]|\\.)*"""") to 0xFFFF8F00.toInt(), // string
      Regex("""\b(fun|val|var|class|data|object|if|else|when|return|
                  |for|while|import|package|override|suspend|null|true|false)\b""".trimMargin("|"))
                              to 0xFF0097A7.toInt(),       // keyword
      Regex("""\b[A-Z][A-Za-z0-9]+\b""") to 0xFF1565C0.toInt(), // type
      Regex("""\b\d+\.?\d*[fLd]?\b""")   to 0xFF8D6E63.toInt(), // number
    )
    Language.PYTHON -> listOf(
      Regex("""#.*$""")              to 0xFF78909C.toInt(),
      Regex(""""[^"]*"|'[^']*'""")   to 0xFFFF8F00.toInt(),
      Regex("""\b(def|class|import|from|if|elif|else|for|while|return|
                  |True|False|None|and|or|not|in|is|lambda)\b""".trimMargin("|"))
                              to 0xFF0097A7.toInt(),
      Regex("""\b\d+\.?\d*\b""")     to 0xFF8D6E63.toInt(),
    )
    Language.SQL -> listOf(
      Regex("""--.*$""")             to 0xFF78909C.toInt(),
      Regex(""""[^"]*"|'[^']*'""")   to 0xFFFF8F00.toInt(),
      Regex("""\b(SELECT|FROM|WHERE|JOIN|LEFT|RIGHT|INNER|ON|AS|GROUP|BY|ORDER|
                  |HAVING|INSERT|UPDATE|DELETE|CREATE|TABLE|INDEX|DROP|ALTER|
                  |AND|OR|NOT|NULL|COUNT|SUM|AVG|MAX|MIN|DISTINCT|LIMIT|OFFSET)\b""".trimMargin("|"),
             RegexOption.IGNORE_CASE) to 0xFF0097A7.toInt(),
      Regex("""\b\d+\.?\d*\b""")     to 0xFF8D6E63.toInt(),
    )
    Language.HTML, Language.XML -> listOf(
      Regex("""<!--.*?-->""", RegexOption.DOT_MATCHES_ALL) to 0xFF78909C.toInt(),
      Regex("""<[^>]+>""")           to 0xFF1565C0.toInt(),  // tag
      Regex(""""[^"]*"""")           to 0xFFFF8F00.toInt(),  // attribute value
    )
    else -> listOf(Regex(""".*""") to 0xFF212121.toInt())
  }

  fun highlight(line: String): List<TextSpan> {
    if (line.isEmpty()) return emptyList()
    // Find all pattern matches, sort by position, fill gaps with default colour
    data class Match(val start: Int, val end: Int, val color: Int)
    val matches = mutableListOf<Match>()

    patterns.forEach { (regex, color) ->
      regex.findAll(line).forEach { result ->
        matches += Match(result.range.first, result.range.last + 1, color)
      }
    }
    // Sort and de-overlap — first match wins at each position
    matches.sortBy { it.start }
    val spans   = mutableListOf<TextSpan>()
    var cursor  = 0
    matches.forEach { m ->
      if (m.start > cursor) {
        spans += TextSpan(line.substring(cursor, m.start), 0xFF212121.toInt())
      }
      if (m.start >= cursor) {
        spans += TextSpan(line.substring(m.start, m.end), m.color)
        cursor = m.end
      }
    }
    if (cursor < line.length) {
      spans += TextSpan(line.substring(cursor), 0xFF212121.toInt())
    }
    return spans
  }
}

enum class Language {
  KOTLIN, JAVA, PYTHON, JAVASCRIPT, TYPESCRIPT, DART,
  SQL, HTML, XML, CSS, JSON, YAML, MARKDOWN,
  C, CPP, RUST, GO, BASH, PHP, RUBY, SWIFT, PLAINTEXT
}
\end{lstlisting}

\section{Text/Code Performance Targets}

\begin{center}\ra{1.35}
\begin{tabular}{L{4.4cm} C{2.8cm} C{2.2cm} L{4.4cm}}
\TH{File} & \TH{QuickEdit / Acode} & \TH{FLUX Target} & \TH{Why FLUX Wins} \\
\rowcolor{rowA} 1\,MB Kotlin file (30k lines) & 3--8\,s to open, slow scroll & \textbf{$<$30ms, 60fps} & LineIndex = 120KB, only visible lines drawn \\
\rowcolor{rowB} 10\,MB log file (300k lines) & ANR / crash & \textbf{$<$50ms, 60fps} & LineIndex = 1.2MB, mmap source \\
\rowcolor{rowA} 50k-line SQL file & Crash or 30+s & \textbf{$<$40ms, 60fps} & SQL syntax highlight, virtual render \\
\rowcolor{rowB} Search in 10MB file & 10--30\,s & \textbf{$<$500ms} & mmap sequential scan at \textasciitilde 2\,GB/s \\
\rowcolor{rowA} Edit + save & 1--3\,s re-render & \textbf{$<$16ms} & Re-layout only modified lines \\
\end{tabular}
\end{center}

\PB

% ============================================================
% CHAPTER 8 — JSON / XML / YAML ENGINE
% ============================================================
\chapter{JSON / XML / YAML Engine — Streaming Tree Viewer}

\section{The Problem With Large Structured Files}

A 50\,MB JSON file with 500,000 nodes cannot be parsed into a tree in RAM
— it would take 500\,MB+ (each JSON node as a Java object is 48--120 bytes
overhead). The solution is \textbf{streaming parse + lazy expansion}: parse
only the top-level structure, expand each subtree on user tap.

\begin{lstlisting}[language=Kotlin,caption={Streaming JSON tree — lazy expansion}]
// A node in the lazy tree
sealed class JsonNode {
  data class JsonString (val value:  String)          : JsonNode()
  data class JsonNumber (val value:  Double)          : JsonNode()
  data class JsonBoolean(val value:  Boolean)         : JsonNode()
  object JsonNull                                     : JsonNode()
  data class JsonObject (
    val keys:     List<String>,
    val offset:   Long,    // byte offset in file where this object starts
    val expanded: Boolean = false,
    val children: Map<String, JsonNode>? = null  // null = not yet loaded
  ) : JsonNode()
  data class JsonArray  (
    val size:     Int,
    val offset:   Long,
    val expanded: Boolean = false,
    val children: List<JsonNode>? = null
  ) : JsonNode()
}

class JsonStreamParser(private val source: MmapSource) {
  // Parse only the top-level structure (keys of root object)
  // Child objects are stored as byte-offset references, not parsed
  fun parseTopLevel(): JsonNode {
    var pos = skipWhitespace(0)
    return when (source.readByte(pos).toInt().toChar()) {
      '{' -> parseObjectShallow(pos)
      '[' -> parseArrayShallow(pos)
      else -> parseScalar(pos).first
    }
  }

  // Shallow parse: read only the keys, store child offsets
  private fun parseObjectShallow(startPos: Long): JsonNode.JsonObject {
    val keys = mutableListOf<String>()
    var pos  = startPos + 1  // skip '{'
    while (pos < source.size) {
      pos = skipWhitespace(pos)
      if (source.readByte(pos).toInt().toChar() == '}') break
      val (key, afterKey) = parseString(pos)
      keys.add(key)
      pos = afterKey
      pos = skipTo(pos, ':') + 1  // skip ':'
      pos = skipValue(pos)         // SKIP the value — don't parse it
      pos = skipWhitespace(pos)
      if (source.readByte(pos).toInt().toChar() == ',') pos++
    }
    return JsonNode.JsonObject(keys, startPos, expanded = false, children = null)
  }

  // Expand a node on user tap — parse only THAT object's children
  fun expandNode(node: JsonNode.JsonObject): JsonNode.JsonObject {
    val children = mutableMapOf<String, JsonNode>()
    var pos = node.offset + 1  // skip '{'
    while (pos < source.size) {
      pos = skipWhitespace(pos)
      if (source.readByte(pos).toInt().toChar() == '}') break
      val (key, afterKey) = parseString(pos)
      pos = afterKey
      pos = skipTo(pos, ':') + 1
      pos = skipWhitespace(pos)
      val (value, afterValue) = parseValue(pos)
      children[key] = value
      pos = afterValue
      pos = skipWhitespace(pos)
      if (source.readByte(pos).toInt().toChar() == ',') pos++
    }
    return node.copy(expanded = true, children = children)
  }

  // Skip a JSON value without parsing it — O(length of value)
  private fun skipValue(pos: Long): Long { /* ... */ return pos }
  private fun skipWhitespace(pos: Long): Long { /* ... */ return pos }
  private fun parseString(pos: Long): Pair<String, Long> { /* ... */ return Pair("", pos) }
  private fun parseValue(pos: Long): Pair<JsonNode, Long> { /* ... */ return Pair(JsonNull, pos) }
  private fun skipTo(pos: Long, ch: Char): Long { /* ... */ return pos }
}
\end{lstlisting}

\section{JSON / XML / YAML Performance Targets}

\begin{center}\ra{1.35}
\begin{tabular}{L{3.6cm} C{2.8cm} C{2.0cm} L{5.4cm}}
\TH{File} & \TH{Current (text editor)} & \TH{FLUX Target} & \TH{How} \\
\rowcolor{rowA} 1\,MB JSON (10k nodes) & 2--5\,s, OOM risk & \textbf{$<$50ms} & Shallow parse top level, lazy expand \\
\rowcolor{rowB} 50\,MB JSON & Crash / OOM & \textbf{$<$100ms} & mmap + top-level only (no full parse) \\
\rowcolor{rowA} 100\,MB XML & Crash / OOM & \textbf{$<$150ms} & SAX streaming, no DOM \\
\rowcolor{rowB} 5\,MB YAML & 3--8\,s & \textbf{$<$60ms} & Custom streaming parser \\
\rowcolor{rowA} Expand node (tap) & N/A (all loaded) & \textbf{$<$30ms} & Parse that subtree on demand from mmap \\
\end{tabular}
\end{center}

\PB

% ============================================================
% CHAPTER 9 — VIDEO / AUDIO ENGINE
% ============================================================
\chapter{Video \& Audio Engine — Native MediaPlayer / ExoPlayer}

\section{Video: Why We Use MediaPlayer (Not ExoPlayer) for Local Files}

\texttt{android.media.MediaPlayer} is the Android OS's built-in video player.
It is hardware-accelerated by the device's video decoder chip (Qualcomm Adreno,
ARM Mali). For local files, it is the \emph{fastest possible} option — it uses
the same decoder as the camera app and Google Photos.

\textbf{ExoPlayer} (now \texttt{androidx.media3}) adds 2\,MB AAR + 500ms
initialisation overhead. It is designed for streaming (HLS, DASH, DRM). For
local MP4/MKV/AVI files, it offers zero benefit over \texttt{MediaPlayer}.

\begin{wbox}[When to Use ExoPlayer vs MediaPlayer]
\begin{itemize}[leftmargin=1.4em,topsep=2pt,itemsep=3pt]
\item \textbf{Use MediaPlayer:} MP4, MKV, AVI, MOV, 3GP, WebM local files.
  Zero library dependency.
\item \textbf{Use ExoPlayer:} HLS streams, DASH streams, DRM-protected content,
  custom data sources. Only if user opens a URL, not a local file.
\end{itemize}
\end{wbox}

\begin{lstlisting}[language=Kotlin,caption={FluxVideoPlayer — minimal local video player}]
class FluxVideoPlayer(
  private val context:     Context,
  private val surfaceView: SurfaceView,
) {
  private var player:   MediaPlayer? = null
  private var prepared  = false

  data class VideoInfo(
    val durationMs: Long,
    val width:      Int,
    val height:     Int,
    val mimeType:   String,
  )

  fun open(file: File, onReady: (VideoInfo) -> Unit) {
    release()
    player = MediaPlayer().apply {
      // Set data source — OS handles file:// URI natively
      setDataSource(context, Uri.fromFile(file))

      // Render directly to SurfaceView — hardware decoder path
      setDisplay(surfaceView.holder)

      // Async prepare — returns immediately, calls back when ready
      setOnPreparedListener { mp ->
        prepared = true
        onReady(VideoInfo(
          durationMs = mp.duration.toLong(),
          width      = mp.videoWidth,
          height     = mp.videoHeight,
          mimeType   = "video/*",
        ))
      }
      setOnErrorListener { _, what, extra ->
        Log.e("FLUX", "MediaPlayer error: what=$what extra=$extra")
        false
      }
      prepareAsync()  // non-blocking: typically ready in 50--200ms for local files
    }
  }

  fun play()  = player?.start()
  fun pause() = player?.pause()
  fun seekTo(positionMs: Long) = player?.seekTo(positionMs.toInt())
  fun release() { player?.release(); player = null; prepared = false }

  val currentPositionMs: Long get() = player?.currentPosition?.toLong() ?: 0
  val isPlaying: Boolean      get() = player?.isPlaying ?: false
}
\end{lstlisting}

\section{Video Thumbnail Generation — Native ThumbnailUtils}

\begin{lstlisting}[language=Kotlin,caption={Zero-library video thumbnail — one API call}]
// API 29+: hardware-accelerated, crops to square, returns RGB_565
suspend fun generateVideoThumb(file: File, size: Size): Bitmap =
  withContext(Dispatchers.IO) {
    ThumbnailUtils.createVideoThumbnail(file, size, null)
    // size = Size(256, 256): renders a 256x256 keyframe thumbnail
    // 'null' cancellationSignal: no timeout
  }

// API 21-28 fallback:
suspend fun generateVideoThumbLegacy(file: File): Bitmap? =
  withContext(Dispatchers.IO) {
    @Suppress("DEPRECATION")
    ThumbnailUtils.createVideoThumbnail(
      file.absolutePath,
      MediaStore.Video.Thumbnails.MINI_KIND  // 512x384
    )
  }
\end{lstlisting}

\section{Audio: MediaPlayer + Waveform Visualiser}

\begin{lstlisting}[language=Kotlin,caption={FluxAudioPlayer with waveform extraction}]
class FluxAudioPlayer(private val file: File) {
  private val player = MediaPlayer().apply {
    setDataSource(file.absolutePath)
    prepare()  // synchronous OK for audio: typically <50ms
  }

  // Extract waveform samples for visualisation
  // MediaExtractor + AudioTrack: reads raw PCM without playing
  suspend fun extractWaveform(samples: Int = 500): FloatArray =
    withContext(Dispatchers.Default) {
      val extractor = MediaExtractor()
      extractor.setDataSource(file.absolutePath)

      // Find audio track
      var audioTrack = -1
      for (i in 0 until extractor.trackCount) {
        val format = extractor.getTrackFormat(i)
        if (format.getString(MediaFormat.KEY_MIME)?.startsWith("audio/") == true) {
          audioTrack = i; break
        }
      }
      if (audioTrack < 0) return@withContext FloatArray(samples)
      extractor.selectTrack(audioTrack)

      val format     = extractor.getTrackFormat(audioTrack)
      val duration   = format.getLong(MediaFormat.KEY_DURATION) // microseconds
      val sampleStep = duration / samples

      val waveform   = FloatArray(samples)
      val buf        = ByteBuffer.allocate(65536)

      for (i in 0 until samples) {
        extractor.seekTo(i * sampleStep, MediaExtractor.SEEK_TO_CLOSEST_SYNC)
        val bytesRead = extractor.readSampleData(buf, 0)
        if (bytesRead > 0) {
          // RMS of this chunk = amplitude at sample point
          var sumSq = 0.0
          for (b in 0 until minOf(bytesRead, 512)) {
            val s = buf.get(b).toFloat() / 128f
            sumSq += s * s
          }
          waveform[i] = Math.sqrt(sumSq / minOf(bytesRead, 512)).toFloat()
        }
      }
      extractor.release()
      waveform
    }
}
\end{lstlisting}

\PB

% ============================================================
% CHAPTER 10 — ZIP / ARCHIVE ENGINE
% ============================================================
\chapter{Archive Engine — ZIP, RAR, 7Z, APK}

\section{ZIP — Built-In, Zero Library}

Android's \texttt{java.util.zip.ZipFile} (built into the JVM) handles ZIP
archives natively. Reading the central directory (list of all files) is O(1)
— it's stored at the end of the file and indexed immediately.

\begin{lstlisting}[language=Kotlin,caption={FluxArchiveReader — tree view of ZIP contents}]
data class ArchiveEntry(
  val path:        String,
  val size:        Long,
  val compressedSize: Long,
  val isDirectory: Boolean,
  val mtime:       Long,
  val children:    MutableList<ArchiveEntry> = mutableListOf(),
)

class FluxZipReader(file: File) : AutoCloseable {
  private val zip = ZipFile(file)

  // Build tree in O(n entries) — typically < 10ms even for large ZIPs
  fun readTree(): ArchiveEntry {
    val root     = ArchiveEntry("/", 0, 0, true, 0)
    val dirMap   = mutableMapOf<String, ArchiveEntry>("/" to root)

    zip.entries().asSequence().forEach { entry ->
      val parts    = entry.name.trimEnd('/').split('/')
      val parent   = ensureDir(parts.dropLast(1), dirMap, root)
      val node     = ArchiveEntry(
        path           = parts.last(),
        size           = entry.size,
        compressedSize = entry.compressedSize,
        isDirectory    = entry.isDirectory,
        mtime          = entry.time,
      )
      parent.children.add(node)
      if (entry.isDirectory) dirMap[entry.name.trimEnd('/')] = node
    }
    return root
  }

  // Extract single entry to destination — streaming, no temp full-file copy
  suspend fun extractEntry(entryName: String, dest: File) =
    withContext(Dispatchers.IO) {
      val entry = zip.getEntry(entryName)
        ?: error("Entry not found: $entryName")
      dest.parentFile?.mkdirs()
      zip.getInputStream(entry).use { input ->
        dest.outputStream().use { output ->
          input.copyTo(output, bufferSize = 65536)
        }
      }
    }

  // Read entry as text (for code preview inside ZIP)
  suspend fun readEntryText(entryName: String, maxBytes: Int = 65536): String =
    withContext(Dispatchers.IO) {
      val entry = zip.getEntry(entryName) ?: return@withContext ""
      zip.getInputStream(entry).use { stream ->
        val buf = ByteArray(minOf(maxBytes.toLong(), entry.size).toInt())
        stream.read(buf)
        String(buf, Charsets.UTF_8)
      }
    }

  override fun close() = zip.close()

  private fun ensureDir(
    parts:  List<String>,
    dirMap: MutableMap<String, ArchiveEntry>,
    root:   ArchiveEntry,
  ): ArchiveEntry {
    if (parts.isEmpty()) return root
    val path = parts.joinToString("/")
    dirMap[path]?.let { return it }
    val parent = ensureDir(parts.dropLast(1), dirMap, root)
    val dir    = ArchiveEntry(parts.last(), 0, 0, true, 0)
    parent.children.add(dir)
    dirMap[path] = dir
    return dir
  }
}
\end{lstlisting}

\section{APK — ZIP + Manifest Parser}

An APK is a ZIP file with a compiled binary manifest
(\texttt{AndroidManifest.xml}) in the AXML (Android Binary XML) format.
We parse this to show app metadata without installing.

\begin{lstlisting}[language=Kotlin,caption={ApkInfoReader — no library needed}]
data class ApkInfo(
  val packageName:  String,
  val versionName:  String,
  val versionCode:  Long,
  val minSdk:       Int,
  val targetSdk:    Int,
  val icon:         Bitmap?,
  val permissions:  List<String>,
  val activities:   List<String>,
  val fileSizeBytes: Long,
)

class ApkInfoReader(private val apkFile: File) {
  fun read(): ApkInfo {
    // PackageManager can parse APK metadata without installing
    val pm   = context.packageManager
    val info = pm.getPackageArchiveInfo(
      apkFile.absolutePath,
      PackageManager.GET_PERMISSIONS or PackageManager.GET_ACTIVITIES
    ) ?: error("Cannot parse APK")

    // Load icon from APK resources
    info.applicationInfo.sourceDir     = apkFile.absolutePath
    info.applicationInfo.publicSourceDir = apkFile.absolutePath
    val icon = try {
      val drawable = info.applicationInfo.loadIcon(pm)
      (drawable as? BitmapDrawable)?.bitmap
        ?: Bitmap.createBitmap(48, 48, Bitmap.Config.ARGB_8888).also {
           Canvas(it).drawBitmap((drawable as? BitmapDrawable)?.bitmap
             ?: createBitmap(48, 48), 0f, 0f, null)
         }
    } catch (e: Exception) { null }

    return ApkInfo(
      packageName   = info.packageName ?: "",
      versionName   = info.versionName ?: "unknown",
      versionCode   = PackageInfoCompat.getLongVersionCode(info),
      minSdk        = info.applicationInfo.minSdkVersion,
      targetSdk     = info.applicationInfo.targetSdkVersion,
      icon          = icon,
      permissions   = info.requestedPermissions?.toList() ?: emptyList(),
      activities    = info.activities?.map { it.name } ?: emptyList(),
      fileSizeBytes = apkFile.length(),
    )
  }
}
\end{lstlisting}

\PB

% ============================================================
% CHAPTER 11 — SVG / FONT / EPUB ENGINE
% ============================================================
\chapter{SVG, Font Preview, EPUB, CSV, Markdown Engines}

\section{SVG — Canvas Renderer}

Android's \texttt{android.graphics.drawable.PictureDrawable} supports SVG
via the \texttt{AndroidSVG} library — but that is 450KB of external code.
For the most common SVG subset (shapes, paths, fills, basic transforms),
we parse directly to \texttt{Path} objects and draw on Canvas.

\begin{lstlisting}[language=Kotlin,caption={Minimal SVG Canvas renderer}]
class FluxSvgRenderer(private val source: MmapSource) {
  data class SvgPath(val path: Path, val fill: Int, val stroke: Int,
                     val strokeWidth: Float)
  private val paths = mutableListOf<SvgPath>()
  var viewBox = RectF(0f, 0f, 100f, 100f)

  fun parse() {
    val parser = XmlPullParserFactory.newInstance().newPullParser()
    parser.setInput(source.readBytes(0, source.size.toInt())
      .inputStream(), "UTF-8")

    val paint = SvgPaintState()
    var event = parser.next()
    while (event != XmlPullParser.END_DOCUMENT) {
      if (event == XmlPullParser.START_TAG) {
        val fill   = parseSvgColor(parser.getAttributeValue(null,"fill")   ?: paint.fill)
        val stroke = parseSvgColor(parser.getAttributeValue(null,"stroke") ?: "none")
        val sw     = parser.getAttributeValue(null,"stroke-width")?.toFloatOrNull() ?: 1f

        when (parser.name) {
          "svg" -> {
            parser.getAttributeValue(null,"viewBox")?.split(" ")?.let { vb ->
              if (vb.size == 4)
                viewBox = RectF(vb[0].toFloat(), vb[1].toFloat(),
                                vb[2].toFloat(), vb[3].toFloat())
            }
          }
          "rect" -> {
            val x = parser.getAttributeValue(null,"x")?.toFloatOrNull() ?: 0f
            val y = parser.getAttributeValue(null,"y")?.toFloatOrNull() ?: 0f
            val w = parser.getAttributeValue(null,"width")?.toFloatOrNull() ?: 0f
            val h = parser.getAttributeValue(null,"height")?.toFloatOrNull() ?: 0f
            val rx= parser.getAttributeValue(null,"rx")?.toFloatOrNull() ?: 0f
            paths += SvgPath(Path().apply { addRoundRect(
              RectF(x,y,x+w,y+h), rx, rx, Path.Direction.CW) },
              fill, stroke, sw)
          }
          "circle" -> {
            val cx = parser.getAttributeValue(null,"cx")?.toFloatOrNull() ?: 0f
            val cy = parser.getAttributeValue(null,"cy")?.toFloatOrNull() ?: 0f
            val r  = parser.getAttributeValue(null,"r")?.toFloatOrNull() ?: 0f
            paths += SvgPath(Path().apply { addCircle(cx,cy,r,Path.Direction.CW) },
              fill, stroke, sw)
          }
          "path" -> {
            val d = parser.getAttributeValue(null,"d") ?: ""
            PathParser.createPathFromPathData(d)?.let { p ->
              paths += SvgPath(p, fill, stroke, sw)
            }
          }
        }
      }
      event = parser.next()
    }
  }

  fun draw(canvas: Canvas, displayRect: RectF) {
    val scaleX = displayRect.width()  / viewBox.width()
    val scaleY = displayRect.height() / viewBox.height()
    val matrix = Matrix().apply {
      setScale(scaleX, scaleY)
      postTranslate(displayRect.left - viewBox.left * scaleX,
                    displayRect.top  - viewBox.top  * scaleY)
    }
    val paint = Paint().apply { isAntiAlias = true }
    paths.forEach { sp ->
      val mapped = Path().apply { sp.path.transform(matrix, this) }
      if (sp.fill != 0) {
        paint.style = Paint.Style.FILL
        paint.color = sp.fill
        canvas.drawPath(mapped, paint)
      }
      if (sp.stroke != 0) {
        paint.style = Paint.Style.STROKE
        paint.color = sp.stroke
        paint.strokeWidth = sp.strokeWidth * minOf(scaleX, scaleY)
        canvas.drawPath(mapped, paint)
      }
    }
  }
}
\end{lstlisting}

\section{CSV — Virtual Column Grid}

CSV files share the same virtual grid architecture as XLSX, but simpler: no
styles, no formulas. We parse lazily using mmap:

\begin{lstlisting}[language=Kotlin,caption={CsvParser — streaming, zero full-load}]
class CsvParser(private val source: MmapSource) {
  // Line offsets built from mmap — same as LineIndex in text renderer
  private val lineOffsets = buildLineOffsets()

  val rowCount: Int get() = lineOffsets.size

  // Parse single row on demand — O(row length)
  fun getRow(rowIndex: Int): List<String> {
    val offset = lineOffsets[rowIndex]
    val end    = if (rowIndex + 1 < lineOffsets.size)
                   lineOffsets[rowIndex + 1] else source.size
    val line   = String(source.readBytes(offset, (end - offset).toInt()))
    return parseCsvLine(line)
  }

  // RFC 4180 compliant parser: handles quoted fields with commas + newlines
  private fun parseCsvLine(line: String): List<String> {
    val fields = mutableListOf<String>()
    var i = 0; val sb = StringBuilder()
    while (i < line.length) {
      when {
        line[i] == '"' -> {  // quoted field
          i++
          while (i < line.length) {
            if (line[i] == '"' && i+1 < line.length && line[i+1] == '"') {
              sb.append('"'); i += 2
            } else if (line[i] == '"') { i++; break }
            else { sb.append(line[i++]) }
          }
        }
        line[i] == ',' -> { fields += sb.toString(); sb.clear(); i++ }
        else -> sb.append(line[i++])
      }
    }
    fields += sb.toString()
    return fields
  }
}
\end{lstlisting}

\section{Markdown — AST + Canvas Render}

\begin{lstlisting}[language=Kotlin,caption={Minimal Markdown renderer to Canvas}]
sealed class MdBlock {
  data class Heading(val level: Int, val text: String)  : MdBlock()
  data class Paragraph(val spans: List<MdSpan>)         : MdBlock()
  data class CodeBlock(val lang: String, val code: String) : MdBlock()
  data class BulletItem(val depth: Int, val spans: List<MdSpan>) : MdBlock()
  data class HorizontalRule(val dummy: Unit = Unit)     : MdBlock()
  data class BlockQuote(val inner: List<MdBlock>)       : MdBlock()
}

sealed class MdSpan {
  data class Plain (val text: String) : MdSpan()
  data class Bold  (val text: String) : MdSpan()
  data class Italic(val text: String) : MdSpan()
  data class Code  (val text: String) : MdSpan()
  data class Link  (val text: String, val url: String) : MdSpan()
}

object MarkdownParser {
  fun parse(source: String): List<MdBlock> {
    val blocks = mutableListOf<MdBlock>()
    val lines  = source.lines()
    var i = 0
    while (i < lines.size) {
      val line = lines[i]
      when {
        line.startsWith("######") -> blocks += MdBlock.Heading(6, line.drop(6).trim())
        line.startsWith("#####")  -> blocks += MdBlock.Heading(5, line.drop(5).trim())
        line.startsWith("####")   -> blocks += MdBlock.Heading(4, line.drop(4).trim())
        line.startsWith("###")    -> blocks += MdBlock.Heading(3, line.drop(3).trim())
        line.startsWith("##")     -> blocks += MdBlock.Heading(2, line.drop(2).trim())
        line.startsWith("# ")     -> blocks += MdBlock.Heading(1, line.drop(2).trim())
        line.startsWith("```")    -> {
          val lang = line.drop(3).trim()
          val code = StringBuilder()
          i++
          while (i < lines.size && !lines[i].startsWith("```")) {
            code.appendLine(lines[i++])
          }
          blocks += MdBlock.CodeBlock(lang, code.toString())
        }
        line.startsWith("- ") || line.startsWith("* ") ->
          blocks += MdBlock.BulletItem(0, parseInline(line.drop(2)))
        line.matches(Regex("""[-*_]{3,}""")) ->
          blocks += MdBlock.HorizontalRule()
        line.isNotBlank() ->
          blocks += MdBlock.Paragraph(parseInline(line))
      }
      i++
    }
    return blocks
  }

  private fun parseInline(text: String): List<MdSpan> {
    // Regex-based inline parsing: **bold**, *italic*, `code`, [link](url)
    val spans   = mutableListOf<MdSpan>()
    val pattern = Regex("""\*\*(.+?)\*\*|\*(.+?)\*|`(.+?)`|\[(.+?)\]\((.+?)\)""")
    var last    = 0
    pattern.findAll(text).forEach { m ->
      if (m.range.first > last) spans += MdSpan.Plain(text.substring(last, m.range.first))
      when {
        m.groups[1] != null -> spans += MdSpan.Bold(m.groups[1]!!.value)
        m.groups[2] != null -> spans += MdSpan.Italic(m.groups[2]!!.value)
        m.groups[3] != null -> spans += MdSpan.Code(m.groups[3]!!.value)
        m.groups[4] != null -> spans += MdSpan.Link(m.groups[4]!!.value, m.groups[5]!!.value)
      }
      last = m.range.last + 1
    }
    if (last < text.length) spans += MdSpan.Plain(text.substring(last))
    return spans
  }
}
\end{lstlisting}

\section{Font Preview}

\begin{lstlisting}[language=Kotlin,caption={Font preview — Typeface.createFromFile()}]
class FontPreviewer(private val file: File) {
  // Android can load TTF, OTF, WOFF directly
  private val typeface: Typeface? = try {
    Typeface.createFromFile(file)  // built-in API, zero library
  } catch (e: Exception) { null }

  fun drawPreview(canvas: Canvas, width: Int, height: Int) {
    typeface ?: return
    val paint = Paint().apply {
      isAntiAlias = true
      this.typeface = typeface
    }
    val preview = listOf(
      Pair(48f, "Aa Bb Cc"),
      Pair(28f, "ABCDEFGHIJKLMNOPQRSTUVWXYZ"),
      Pair(24f, "abcdefghijklmnopqrstuvwxyz"),
      Pair(22f, "0123456789 !@#\$%^&*()"),
      Pair(18f, "The quick brown fox jumps over the lazy dog."),
    )
    var y = 60f
    preview.forEach { (size, text) ->
      paint.textSize = size
      paint.color    = 0xFF212121.toInt()
      canvas.drawText(text, 24f, y, paint)
      y += size * 1.6f
    }
  }
}
\end{lstlisting}

\PB

% ============================================================
% CHAPTER 12 — HTML ENGINE
% ============================================================
\chapter{HTML Engine — WebView Optimised for Local Files}

\section{WebView Is Already Loaded}

Android's \texttt{WebView} is backed by the system WebView provider
(typically Chrome or Android System WebView). It is already loaded into
memory as a system process. The cold start cost is \textbf{not} from loading
the rendering engine — it is from creating the \texttt{WebView} instance in
your process and waiting for the first \texttt{onPageFinished} callback.

\begin{lstlisting}[language=Kotlin,caption={FluxWebViewer — optimised for local HTML files}]
class FluxWebViewer(context: Context) : WebView(context) {
  init {
    settings.apply {
      // Local file access — needed for HTML with local images/CSS
      allowFileAccess           = true
      allowContentAccess        = true

      // Performance
      cacheMode                 = WebSettings.LOAD_CACHE_ELSE_NETWORK
      setRenderPriority(WebSettings.RenderPriority.HIGH)
      blockNetworkLoads         = true  // offline-only: no network requests
      loadsImagesAutomatically  = true
      domStorageEnabled         = false // not needed for static HTML
      javaScriptEnabled         = false // disable by default; enable if needed
      useWideViewPort           = true
      loadWithOverviewMode      = true
      setSupportZoom(true)
      builtInZoomControls       = true
      displayZoomControls       = false
    }

    // Intercept any network requests (should be none for local files)
    webViewClient = object : WebViewClient() {
      override fun shouldInterceptRequest(
        view: WebView, request: WebResourceRequest
      ): WebResourceResponse? {
        // Block all network requests for privacy + performance
        return if (request.url.scheme != "file") {
          WebResourceResponse("text/plain", "UTF-8",
            "blocked".byteInputStream())
        } else null
      }

      override fun onPageFinished(view: WebView, url: String) {
        // Inject minimal CSS fixes for better mobile rendering
        view.evaluateJavascript("""
          (function() {
            var style = document.createElement('style');
            style.textContent = 'body { max-width: 100%; word-wrap: break-word; }
              img { max-width: 100%; height: auto; }
              pre { overflow-x: auto; }';
            document.head.appendChild(style);
          })();
        """.trimIndent(), null)
      }
    }
  }

  fun openLocalFile(file: File) {
    // loadUrl with file:// is fastest for local HTML
    loadUrl("file://${file.absolutePath}")
  }

  fun openHtmlString(html: String, baseUrl: String? = null) {
    loadDataWithBaseURL(baseUrl, html, "text/html", "UTF-8", null)
  }
}
\end{lstlisting}

\section{EPUB — HTML + ZIP}

An EPUB is a ZIP containing HTML files + a manifest. We parse the manifest
to get reading order, then render each HTML chapter in WebView.

\begin{lstlisting}[language=Kotlin,caption={EpubReader — manifest + chapter ordering}]
data class EpubChapter(val title: String, val href: String)

class EpubReader(private val file: File) {
  fun readManifest(): List<EpubChapter> {
    val chapters = mutableListOf<EpubChapter>()
    ZipFile(file).use { zip ->
      // Find OPF file from META-INF/container.xml
      val container = zip.getEntry("META-INF/container.xml") ?: return emptyList()
      val opfPath   = parseContainerXml(zip.getInputStream(container))

      // Parse OPF: spine gives reading order
      val opfEntry  = zip.getEntry(opfPath) ?: return emptyList()
      chapters += parseOpf(zip.getInputStream(opfEntry))
    }
    return chapters
  }

  fun extractChapterToTempDir(chapterHref: String, tempDir: File): File {
    val outFile = File(tempDir, chapterHref.substringAfterLast('/'))
    ZipFile(file).use { zip ->
      zip.getEntry("OEBPS/$chapterHref")?.let { entry ->
        zip.getInputStream(entry).use { it.copyTo(outFile.outputStream()) }
      }
    }
    return outFile
  }
}
\end{lstlisting}

\PB

% ============================================================
% CHAPTER 13 — EDIT LAYER
% ============================================================
\chapter{Edit Layer — Universal In-Memory Edit System}

\section{Design Principle: Edits Are Overlays}

The viewer always reads from the original \emph{immutable} \texttt{MmapSource}.
Edits are stored in a separate \texttt{EditOverlay} that intercepts reads:
when rendering, the engine first checks if the current position has an edit;
if so, it uses the edit data; otherwise it uses the mmap data. On save, edits
are merged and written to a new file (the original is never overwritten until
the write is confirmed complete).

\begin{lstlisting}[language=Kotlin,caption={EditOverlay — unified edit buffer}]
// An edit is a replacement of byte range [start, end) with new bytes
data class Edit(
  val id:       Long,         // monotonically increasing edit ID
  val offset:   Long,        // byte offset in original file
  val length:   Int,          // number of original bytes replaced (0 = insert)
  val newBytes: ByteArray,    // replacement content
  val type:     EditType,
)
enum class EditType { INSERT, REPLACE, DELETE }

class EditOverlay {
  // TreeMap: sorted by offset for fast range queries
  private val edits = TreeMap<Long, Edit>()
  private var nextId = 0L

  // O(log n) insert/replace
  fun applyEdit(offset: Long, length: Int, newBytes: ByteArray): Long {
    val id   = nextId++
    edits[offset] = Edit(id, offset, length, newBytes, when {
      length == 0      -> EditType.INSERT
      newBytes.isEmpty()-> EditType.DELETE
      else             -> EditType.REPLACE
    })
    return id
  }

  fun undoEdit(id: Long) {
    edits.entries.removeIf { it.value.id == id }
  }

  // Read through overlay: applies edits to a virtual byte stream
  fun readVirtual(original: MmapSource, start: Long, length: Int): ByteArray {
    val result   = ByteArrayOutputStream(length)
    var pos      = start
    val end      = start + length

    while (pos < end) {
      // Find next edit at or after pos
      val editEntry = edits.ceilingEntry(pos)
      if (editEntry == null || editEntry.key >= end) {
        // No more edits in this range: read directly from mmap
        val toRead = (end - pos).toInt()
        result.write(original.readBytes(pos, toRead))
        break
      }
      val edit = editEntry.value
      if (edit.offset > pos) {
        // Gap before edit: read from mmap
        result.write(original.readBytes(pos, (edit.offset - pos).toInt()))
      }
      // Apply edit
      result.write(edit.newBytes)
      pos = edit.offset + edit.length
    }
    return result.toByteArray()
  }

  // Save: merge original + edits into new file
  suspend fun saveTo(original: MmapSource, output: File) =
    withContext(Dispatchers.IO) {
      output.parentFile?.mkdirs()
      output.outputStream().buffered(65536).use { out ->
        var pos = 0L
        edits.values.sortedBy { it.offset }.forEach { edit ->
          if (edit.offset > pos) {
            out.write(original.readBytes(pos, (edit.offset - pos).toInt()))
          }
          out.write(edit.newBytes)
          pos = edit.offset + edit.length
        }
        if (pos < original.size) {
          out.write(original.readBytes(pos, (original.size - pos).toInt()))
        }
      }
    }
}
\end{lstlisting}

\section{Format-Specific Edit APIs}

\begin{center}\ra{1.35}
\begin{tabular}{L{2.0cm} L{4.6cm} L{3.6cm} L{3.6cm}}
\TH{Format} & \TH{Edit Operation} & \TH{Edit Overlay Action} & \TH{Save Strategy} \\
\rowcolor{rowA} TXT / Code & Type character at cursor & INSERT edit at caret offset & Stream merge + write \\
\rowcolor{rowB} TXT / Code & Delete selection & DELETE edit (length = selection) & Stream merge + write \\
\rowcolor{rowA} PDF & Add text annotation & EditOverlay stores annotation metadata & Re-render page with annotation layer \\
\rowcolor{rowB} PDF & Highlight region & EditOverlay stores highlight rect + colour & Re-render page with annotation layer \\
\rowcolor{rowA} DOCX & Modify paragraph text & Patch \texttt{word/document.xml} within ZIP & Rewrite XML node + repack ZIP \\
\rowcolor{rowB} XLSX & Edit cell value & Update cell in SheetData + XML & Rewrite \texttt{sheet1.xml} + repack ZIP \\
\rowcolor{rowA} Image & Crop / rotate & New Bitmap op + encode to JPEG & Write new JPEG to temp file \\
\rowcolor{rowB} CSV & Edit cell value & Replace text range in mmap view & Stream merge with new value \\
\end{tabular}
\end{center}

\PB

% ============================================================
% CHAPTER 14 — COMPLETE FORMAT SUPPORT MATRIX
% ============================================================
\chapter{Complete Format Support Matrix}

\section{All Supported File Types}

\begin{center}\ra{1.3}
\begin{longtable}{L{1.8cm} L{2.4cm} L{2.0cm} C{1.4cm} C{1.4cm} L{4.2cm}}
\THG{Category} & \THG{Format} & \THG{Extensions} & \THG{View} & \THG{Edit} &
\THG{Implementation} \\
\endhead
\rowcolor{rowA} Documents & PDF & .pdf & \checkmark & Ann. only & PdfRenderer API + tile cache \\
\rowcolor{rowB} Documents & Word & .docx, .doc & \checkmark & \checkmark & OOXML ZIP+XML parser + Canvas layout \\
\rowcolor{rowA} Documents & Excel & .xlsx, .xls & \checkmark & \checkmark & OOXML ZIP+XML parser + virtual grid \\
\rowcolor{rowB} Documents & PowerPoint & .pptx, .ppt & \checkmark & — & Slide renderer: Canvas shapes + text \\
\rowcolor{rowA} Documents & RTF & .rtf & \checkmark & — & RTF token parser → canvas layout \\
\rowcolor{rowB} Documents & ODT/ODS & .odt, .ods & \checkmark & — & ODF ZIP+XML (same approach as OOXML) \\
\rowcolor{rowA} Images & JPEG/JPG & .jpg, .jpeg & \checkmark & Crop/rotate & BitmapRegionDecoder + tile cache \\
\rowcolor{rowB} Images & PNG & .png & \checkmark & — & BitmapRegionDecoder \\
\rowcolor{rowA} Images & WebP & .webp & \checkmark & — & BitmapRegionDecoder \\
\rowcolor{rowB} Images & HEIC/HEIF & .heic, .heif & \checkmark & — & ImageDecoder API (API 28+) \\
\rowcolor{rowA} Images & GIF & .gif & \checkmark & — & ImageDecoder + AnimatedImageDrawable \\
\rowcolor{rowB} Images & RAW/DNG & .dng, .raw, .cr2 & Thumb+preview & — & ThumbnailUtils thumb, ExifInterface meta \\
\rowcolor{rowA} Images & BMP & .bmp & \checkmark & — & BitmapFactory + inSampleSize \\
\rowcolor{rowB} Images & SVG & .svg & \checkmark & — & Custom Canvas path renderer \\
\rowcolor{rowA} Video & MP4 & .mp4, .m4v & \checkmark & — & MediaPlayer + SurfaceView \\
\rowcolor{rowB} Video & MKV & .mkv & \checkmark & — & MediaPlayer + SurfaceView \\
\rowcolor{rowA} Video & AVI & .avi & \checkmark & — & MediaPlayer + SurfaceView \\
\rowcolor{rowB} Video & MOV & .mov & \checkmark & — & MediaPlayer + SurfaceView \\
\rowcolor{rowA} Video & WebM & .webm & \checkmark & — & MediaPlayer + SurfaceView \\
\rowcolor{rowB} Video & 3GP & .3gp & \checkmark & — & MediaPlayer + SurfaceView \\
\rowcolor{rowA} Audio & MP3 & .mp3 & \checkmark & — & MediaPlayer + waveform visualiser \\
\rowcolor{rowB} Audio & FLAC & .flac & \checkmark & — & MediaPlayer + waveform visualiser \\
\rowcolor{rowA} Audio & AAC/M4A & .aac, .m4a & \checkmark & — & MediaPlayer + waveform visualiser \\
\rowcolor{rowB} Audio & OGG & .ogg & \checkmark & — & MediaPlayer \\
\rowcolor{rowA} Audio & WAV & .wav & \checkmark & — & MediaPlayer + sample-accurate waveform \\
\rowcolor{rowB} Audio & OPUS & .opus & \checkmark & — & MediaPlayer \\
\rowcolor{rowA} Text & Plain text & .txt & \checkmark & \checkmark & Canvas LineRenderer + LineIndex (mmap) \\
\rowcolor{rowB} Text & Markdown & .md & \checkmark & \checkmark & MD parser + Canvas block renderer \\
\rowcolor{rowA} Code & Kotlin/Java & .kt, .java & \checkmark & \checkmark & Text renderer + KT/Java highlighter \\
\rowcolor{rowB} Code & Python & .py & \checkmark & \checkmark & Text renderer + Python highlighter \\
\rowcolor{rowA} Code & JavaScript/TS & .js, .ts & \checkmark & \checkmark & Text renderer + JS highlighter \\
\rowcolor{rowB} Code & Dart & .dart & \checkmark & \checkmark & Text renderer + Dart highlighter \\
\rowcolor{rowA} Code & C/C++/Rust/Go & .c, .cpp, .rs, .go & \checkmark & \checkmark & Text renderer + respective highlighter \\
\rowcolor{rowB} Data & SQL & .sql & \checkmark & \checkmark & Text renderer + SQL highlighter \\
\rowcolor{rowA} Data & JSON & .json & \checkmark & \checkmark & Streaming tree parser + lazy expand \\
\rowcolor{rowB} Data & XML & .xml & \checkmark & \checkmark & SAX streaming + tree viewer \\
\rowcolor{rowA} Data & YAML & .yaml, .yml & \checkmark & \checkmark & Streaming YAML parser \\
\rowcolor{rowB} Data & CSV & .csv, .tsv & \checkmark & \checkmark & LineIndex + virtual grid renderer \\
\rowcolor{rowA} Data & SQLite DB & .db, .sqlite & \checkmark & — & SQLiteDatabase API: table list + row viewer \\
\rowcolor{rowB} Web & HTML & .html, .htm & \checkmark & — & Optimised local WebView \\
\rowcolor{rowA} Web & CSS & .css & \checkmark & \checkmark & Text renderer + CSS highlighter \\
\rowcolor{rowB} Archives & ZIP & .zip & \checkmark & — & ZipFile (built-in) + tree viewer \\
\rowcolor{rowA} Archives & RAR & .rar & View list only & — & RAR5 header parser (minimal, 200 lines) \\
\rowcolor{rowB} Archives & 7Z & .7z & View list only & — & 7Z header parser (minimal) \\
\rowcolor{rowA} Archives & APK & .apk & \checkmark & — & PackageManager + ZipFile \\
\rowcolor{rowB} Archives & EPUB & .epub & \checkmark & — & ZIP + OPF manifest + WebView per chapter \\
\rowcolor{rowA} Fonts & TTF/OTF & .ttf, .otf & \checkmark & — & Typeface.createFromFile() + Canvas preview \\
\rowcolor{rowB} Fonts & WOFF & .woff & \checkmark & — & Typeface.createFromFile() (API 29+) \\
\end{longtable}
\end{center}

\section{Android API Dependency Map}

\begin{center}\ra{1.35}
\begin{tabular}{L{4.0cm} L{3.6cm} C{2.2cm} L{4.0cm}}
\THC{API Used} & \THC{Purpose} & \THC{Min API} & \THC{Formats} \\
\rowcolor{rowA} \texttt{PdfRenderer} & PDF rendering & 21 & PDF \\
\rowcolor{rowB} \texttt{BitmapRegionDecoder} & Image tile decode & 10 & JPEG, PNG, WebP \\
\rowcolor{rowA} \texttt{ImageDecoder} & HEIC, GIF animation & 28 & HEIC, HEIF, GIF, WebP \\
\rowcolor{rowB} \texttt{MediaPlayer} & Video + audio playback & 1 & All video/audio \\
\rowcolor{rowA} \texttt{MediaExtractor} & Waveform extraction & 16 & All audio \\
\rowcolor{rowB} \texttt{ThumbnailUtils} & Video thumbnail & 8 (legacy), 29 (new) & Video \\
\rowcolor{rowA} \texttt{ZipFile} & ZIP archive reading & 1 & ZIP, DOCX, XLSX, APK, EPUB \\
\rowcolor{rowB} \texttt{PackageManager} & APK info parsing & 1 & APK \\
\rowcolor{rowA} \texttt{XmlPullParser} & XML streaming parse & 1 & DOCX, XLSX, PPTX, SVG, HTML \\
\rowcolor{rowB} \texttt{SQLiteDatabase} & SQLite file open & 1 & .db, .sqlite \\
\rowcolor{rowA} \texttt{WebView} & HTML rendering & 1 & HTML, EPUB chapters \\
\rowcolor{rowB} \texttt{Typeface.createFromFile()} & Font rendering & 1 & TTF, OTF, WOFF \\
\rowcolor{rowA} \texttt{PdfDocument} & PDF annotation save & 19 & PDF (edit mode) \\
\rowcolor{rowB} \texttt{StorageManager.openProfd()} & Direct file descriptor & 26 & All (performance) \\
\end{tabular}
\end{center}

\PB

% ============================================================
% CHAPTER 15 — FLUTTER INTEGRATION
% ============================================================
\chapter{Flutter Integration — ViewerRouter \& Platform Views}

\section{ViewerRouter — Format Dispatch}

\begin{lstlisting}[language=Kotlin,caption={ViewerRouter.dart — format-based screen dispatch}]
// Dart — ViewerRouter.dart
class ViewerRouter extends ConsumerWidget {
  final int    fid;
  final String path;
  final String mimeType;

  @override
  Widget build(BuildContext context, WidgetRef ref) {
    return switch(detectFormat(path, mimeType)) {
      FileFormat.PDF     => PdfViewerScreen(fid: fid, path: path),
      FileFormat.DOCX    => DocxViewerScreen(fid: fid, path: path),
      FileFormat.XLSX    => XlsxViewerScreen(fid: fid, path: path),
      FileFormat.PPTX    => PptxViewerScreen(fid: fid, path: path),
      FileFormat.IMAGE   => ImageViewerScreen(fid: fid, path: path),
      FileFormat.VIDEO   => VideoPlayerScreen(fid: fid, path: path),
      FileFormat.AUDIO   => AudioPlayerScreen(fid: fid, path: path),
      FileFormat.TEXT    => TextViewerScreen(fid: fid, path: path),
      FileFormat.CODE    => CodeViewerScreen(fid: fid, path: path,
                             language: detectLanguage(path)),
      FileFormat.JSON    => JsonViewerScreen(fid: fid, path: path),
      FileFormat.XML     => XmlViewerScreen(fid: fid, path: path),
      FileFormat.HTML    => HtmlViewerScreen(fid: fid, path: path),
      FileFormat.ZIP     => ArchiveViewerScreen(fid: fid, path: path),
      FileFormat.APK     => ApkInfoScreen(fid: fid, path: path),
      FileFormat.SVG     => SvgViewerScreen(fid: fid, path: path),
      FileFormat.EPUB    => EpubViewerScreen(fid: fid, path: path),
      FileFormat.CSV     => CsvViewerScreen(fid: fid, path: path),
      FileFormat.SQLITE  => SqliteViewerScreen(fid: fid, path: path),
      FileFormat.FONT    => FontPreviewScreen(fid: fid, path: path),
      _                  => UnknownFileScreen(fid: fid, path: path),
    };
  }
}
\end{lstlisting}

\section{AndroidView — Bridging Native Rendering to Flutter}

For rendering-intensive formats (PDF, images, video), we use
\texttt{AndroidView} or \texttt{PlatformViewLink} to embed the native
Kotlin-drawn view directly inside the Flutter widget tree.

\begin{lstlisting}[language=Kotlin,caption={PdfViewerScreen.dart — embedded native view}]
// Dart — PdfViewerScreen.dart
class PdfViewerScreen extends StatelessWidget {
  final int    fid;
  final String path;

  @override
  Widget build(BuildContext context) {
    return Scaffold(
      body: AndroidView(
        viewType:             'com.flux/pdf_viewer',
        creationParams:       {'path': path, 'fid': fid},
        creationParamsCodec:  const StandardMessageCodec(),
        onPlatformViewCreated: (id) {
          // Connect MethodChannel to this specific view
          _pdfController = PdfViewController(id);
        },
      ),
    );
  }
}

// Kotlin — PdfViewerFactory.kt
class PdfViewerFactory(private val context: Context) :
    PlatformViewFactory(StandardMessageCodec.INSTANCE) {

  override fun create(ctx: Context, viewId: Int, args: Any?): PlatformView {
    val params  = args as Map<String, Any>
    val path    = params["path"] as String
    val view    = FluxPdfView(ctx, File(path))
    return object : PlatformView {
      override fun getView(): View = view
      override fun dispose() { view.release() }
    }
  }
}
\end{lstlisting}

\section{Performance: First Paint Timeline}

\begin{center}\ra{1.35}
\begin{tabular}{L{2.4cm} L{2.0cm} L{2.0cm} L{2.0cm} L{5.4cm}}
\TH{Format} & \TH{T=0ms} & \TH{T=50ms} & \TH{T=100ms} & \TH{What Happens at Each Stage} \\
\rowcolor{rowA} PDF (10 pg) & Route resolved & Tile 0 rendered & Full page visible &
  T=0: FormatDetector. T=30ms: PdfRenderer.openPage(). T=50ms: first 512×512 tile on screen \\
\rowcolor{rowB} DOCX & Route resolved & First para on canvas & Full viewport &
  T=0: ZipFile.open(). T=20ms: document.xml stream starts. T=80ms: first 200 lines laid out \\
\rowcolor{rowA} XLSX & Route resolved & Headers on grid & 20 rows visible &
  T=0: sharedStrings cached. T=30ms: first 50 cells parsed. T=60ms: virtual grid first draw \\
\rowcolor{rowB} Image (12MP) & Thumbnail shown & Low-res (÷8) displayed & Visible region decoded &
  T=0: from ThumbnailCache. T=30ms: BitmapRegionDecoder inSampleSize=8. T=80ms: inSampleSize=1 visible tile \\
\rowcolor{rowA} Text (10MB) & Route resolved & LineIndex built & 30 visible lines drawn &
  T=0: MmapSource open. T=20ms: LineIndex built (scan for \textbackslash n). T=40ms: first screen rendered \\
\rowcolor{rowB} Video & Route resolved & Surface created & First frame displayed &
  T=0: MediaPlayer.setDataSource(). T=30ms: prepareAsync(). T=80ms: onPrepared, start() \\
\end{tabular}
\end{center}

\PB

% ============================================================
% CHAPTER 16 — IMPLEMENTATION PLAN
% ============================================================
\chapter{Implementation Plan — Viewer Engine Build Order}

\section{Build Phases}

\begin{center}\ra{1.4}
\begin{tabular}{C{1.2cm} L{2.2cm} L{2.8cm} L{5.2cm} L{3.4cm}}
\THO{Phase} & \THO{Weeks} & \THO{Name} & \THO{What to Build} & \THO{Acceptance Criteria} \\
\rowcolor{rowA} V1 & 1--2 & Core infra & MmapSource, FormatDetector, TileEngine, EditOverlay &
  Read any file format via mmap. TileEngine renders 512×512 tiles. \\
\rowcolor{rowB} V2 & 3--5 & Image + Video & BitmapRegionDecoder tile renderer, MediaPlayer video, audio player &
  12MP image opens in $<$30ms (thumb). Video plays in $<$80ms. \\
\rowcolor{rowA} V3 & 6--8 & Text + Code & FluxTextRenderer, LineIndex, SyntaxHighlighter (6 languages) &
  10MB file opens in $<$50ms. 60fps scroll on 300k line file. \\
\rowcolor{rowB} V4 & 9--11 & PDF & FluxPdfRenderer tile system, PdfAnnotationLayer &
  10-page PDF first tile $<$80ms. Annotation overlay works. \\
\rowcolor{rowA} V5 & 12--15 & Office formats & DocxParser+layout, XlsxVirtualGrid, PptxSlideRenderer &
  20-page DOCX opens $<$150ms. XLSX 1000 rows: 60fps scroll. \\
\rowcolor{rowB} V6 & 16--18 & Data formats & JSON/XML/YAML tree viewer, CSV virtual grid, SQLite viewer &
  50MB JSON opens $<$100ms. 100k-row CSV: 60fps scroll. \\
\rowcolor{rowA} V7 & 19--21 & Web + Archives & HTML WebView, ZIP/APK viewer, EPUB reader &
  Local HTML opens $<$200ms. ZIP 1000 files tree $<$30ms. \\
\rowcolor{rowB} V8 & 22--24 & Advanced & SVG renderer, Font preview, Markdown, edit + save for all &
  SVG 200 paths renders $<$100ms. Edit+save round-trip verified. \\
\rowcolor{rowA} V9 & 25--26 & Flutter UI & ViewerRouter, all viewer screens, edit toolbar, PlatformView &
  All formats open from FLUX browser with correct viewer. \\
\rowcolor{rowB} V10 & 27 & Polish & Error handling all formats, unsupported fallback, a11y &
  No crash on any file type. Unknown format shows hex viewer. \\
\end{tabular}
\end{center}

\section{AI Coding Agent Rules for Viewer Engine}

\begin{cbox}[FORBIDDEN in Viewer Engine Code]
\begin{enumerate}[leftmargin=1.4em,topsep=2pt,itemsep=4pt]
  \item \textbf{NEVER} call \texttt{File.readBytes()} or \texttt{readText()} for
    files that may exceed 1\,MB. Always use \texttt{MmapSource}.
  \item \textbf{NEVER} create a \texttt{Bitmap} larger than $2048 \times 2048$
    in a single allocation. Always tile.
  \item \textbf{NEVER} use \texttt{Bitmap.Config.ARGB\_8888} for photos/PDF
    tiles. Use \texttt{RGB\_565} (50\% RAM saving, visually identical for content).
  \item \textbf{NEVER} load the full document into a DOM before showing anything.
    Always stream-parse and show first viewport within 150\,ms.
  \item \textbf{NEVER} do rendering (Canvas draws) on the main thread. Tile
    generation always runs on \texttt{Dispatchers.Default}.
  \item \textbf{NEVER} add external rendering libraries for any format already
    achievable with built-in Android APIs.
  \item \textbf{NEVER} overwrite the original file during save. Always write to
    a temp file, then atomic rename.
  \item \textbf{NEVER} enable \texttt{javaScriptEnabled} in WebView unless the
    file explicitly requires it and the user approves.
\end{enumerate}
\end{cbox}

\begin{gbox}[REQUIRED in Every Format Renderer]
\begin{enumerate}[leftmargin=1.4em,topsep=2pt,itemsep=4pt]
  \item Every renderer \textbf{MUST} implement \texttt{AutoCloseable} and
    release all resources in \texttt{close()}.
  \item Every renderer \textbf{MUST} show \emph{something} (thumbnail, placeholder,
    or first tile) within 100\,ms of being opened.
  \item Every renderer \textbf{MUST} handle \texttt{OutOfMemoryError} gracefully:
    reduce tile cache, show error state, never crash.
  \item Every renderer \textbf{MUST} support \texttt{MmapSource} as its input,
    never a raw \texttt{ByteArray} of the full file.
  \item Every tile-based renderer \textbf{MUST} use \texttt{RGB\_565} for
    non-alpha content and \texttt{ARGB\_8888} only when transparency is required.
  \item Every edit operation \textbf{MUST} go through \texttt{EditOverlay},
    never modify the source mmap directly.
  \item Every format's save \textbf{MUST} write to a temp file first, verify
    integrity, then atomic-rename over the original.
\end{enumerate}
\end{gbox}

\vspace{24pt}
\begin{center}
  \color{P1}\rule{0.5\linewidth}{1pt}\\[8pt]
  {\sffamily\bfseries\Large END OF FLUX VIEWER ENGINE DOCUMENT}\\[6pt]
  \rule{0.5\linewidth}{1pt}\\[8pt]
  {\sffamily\small\color{GY}
    Market Analysis · Core Architecture · PDF · DOCX · XLSX · Images\\[3pt]
    Text/Code · JSON/XML · Video/Audio · Archives · SVG · HTML · Edit Layer\\[3pt]
    Complete Format Matrix · Flutter Integration · Implementation Plan\\[4pt]
    Version 1.0 · 2025 · Confidential
  }
\end{center}

\end{document}